% =============================================================================
%  Datei-Hinweise & Konfiguration: siehe config.md
%  -> Abschnitt "main.tex"
%
%  WICHTIG: Vektorgrafiken werden via Paket "svg" (\includesvg) eingebunden.
%  Das erfordert --shell-escape und ein installiertes Inkscape. Lokaler Aufruf:
%      latexmk -pdf -shell-escape main.tex
%  Gesetzt sind: \usepackage{svg}, \svgpath{{abbildungen/}} und
%  \svgsetup{inkscapelatex=false} (bettet SVG-Text fest ins PDF ein und
%  vermeidet ueberlagerte Text-Artefakte in den BPMN-Modellen).
% =============================================================================

\documentclass[
    paper=a4,           % DIN A4 (Leitfaden Kap. 5.2.1)
    fontsize=11pt,      % Open Sans 11 pt analog zu Arial 11 pt
                        % (Leitfaden Kap. 5.2.2: ~37 Zeilen / ~60 Zeichen /
                        % ~2.200 Zeichen je Seite einzuhalten)
    twoside=false,      % einseitiger Druck (Leitfaden Kap. 5.2.1: nur
                        % bei beidseitigem Druck "twoside=true" zusammen
                        % mit BCOR setzen, Kap. 5.10.1)
    open=any,           % Kapitel beginnen auf der naechsten freien Seite,
                        % nicht zwingend auf einer rechten -- vermeidet
                        % unerwuenschte Leerseiten im Oneside-Druck
    BCOR=0mm,           % Bindekorrektur, z. B. 5 mm bei Leimbindung
                        % (Leitfaden Kap. 5.10.1)
    headings=normal,    % normale Ueberschriftsgroessen
                        % (Leitfaden Kap. 5.2.3)
    parskip=half-,      % Absatz: 0 pt vor, 6 pt nach (Leitfaden Kap. 5.3)
    listof=totoc,       % Verzeichnisse im Inhaltsverzeichnis aufgefuehrt
                        % (Leitfaden Kap. 5.1)
    bibliography=totoc, % Quellenverzeichnis im Inhaltsverzeichnis
                        % (Leitfaden Kap. 6.3)
    numbers=noenddot    % keine Schlusspunkte hinter Kapitelnummern
                        % (Leitfaden Kap. 5.2.3)
]{scrreprt}

% ----------------------------- Sprache & Encoding ---------------------------
% Deutsche Hauptsprache mit englischer Zusatzsprache (Leitfaden Kap. 2.2:
% "Anfuehrungszeichen der Dokumentensprache anpassen: Deutsch „...",
% Englisch ''..''."). \foreignlanguage{english}{...} bzw.
% \begin{otherlanguage}{english}...\end{otherlanguage} aktivieren englische
% Trennregeln und Anfuehrungszeichen (sinnvoll fuer in der Informatik
% haeufige englische Fachbegriffe).
\usepackage[utf8]{inputenc}
\usepackage[T1]{fontenc}
\usepackage[main=ngerman,english]{babel}
\usepackage[autostyle=true,german=quotes]{csquotes} % „..." per \enquote{}

% ----------------------------- Schrift --------------------------------------
% FOM-Leitfaden Kap. 5.2.2: "Fuer den Text der Arbeit muss eine Schriftart
% gewaehlt werden, die durchgaengig verwendet wird." Beispiele: Times 12 pt,
% Arial 11 pt; entscheidend ist, dass ca. 37 Zeilen / 60 Zeichen / 2.200
% Zeichen je Seite erreicht werden.
% Open Sans (humanistische Sans-Serif, Apache-2.0-Lizenz) erfuellt diese
% Vorgabe bei 11 pt analog zu Arial und ist daher leitfadenkonform.
\usepackage[default,scale=0.97]{opensans}   % Open Sans als Hauptschrift (auf ~2.200 Zeichen/Seite kalibriert, Leitfaden Kap. 5.2.2)
\usepackage[scaled=0.95]{inconsolata}        % Inconsolata als Monospace-
                                             % schrift fuer Code-Listings
                                             % (FOM Kap. 5.9.2 verlangt eine
                                             % "Fixtype" fuer Quelltexte)
\usepackage{microtype}                       % bessere Mikrotypografie

% ----------------------------- Seitenraender --------------------------------
% Vorgabe Leitfaden Kap. 5.2.1: oben 4 cm, unten 2 cm, links 4 cm, rechts 2 cm,
% Kopfzeile 2 cm, Fusszeile 2 cm Abstand vom Seitenrand.
\usepackage[
    a4paper,
    top=4cm,
    bottom=2cm,
    left=4cm,
    right=2cm,
    headsep=1cm,
    footskip=1cm,
    includeheadfoot=false
]{geometry}

% ----------------------------- Zeilenabstand --------------------------------
% Zeilenabstand 1,5-zeilig im Fliesstext (Leitfaden Kap. 5.2.2:
% "Zeilenabstand im Text: 1,5-zeilig").
\usepackage{setspace}
\onehalfspacing

% ----------------------------- Kopf- und Fusszeile --------------------------
% Kopfzeile mit Abschnittsueberschrift und Seitennummer gemaess
% Leitfaden Kap. 5.2.1 ("Kopfbereich: hier koennen Abschnittsueberschriften
% und Seitennummern erscheinen").
\usepackage{scrlayer-scrpage}
\setlength{\headheight}{28pt}     % Platz fuer zweizeiligen Kopf (Text + Akzent)
\pagestyle{scrheadings}
\clearpairofpagestyles
\automark[chapter]{chapter}       % Hauptabschnitt = Kapitel im Kopf
% Auch Kapitel-Startseiten erhalten das volle Header-Layout (statt
% \pagestyle{plain}), damit der Hauptabschnitts-Titel ausnahmslos sichtbar ist.
\renewcommand*{\chapterpagestyle}{scrheadings}

% Akzentlinien-Breiten entlang einer gemeinsamen Goldener-Schnitt-Familie
% (Hervorhebungen sparsam und einheitlich – Leitfaden Kap. 5.2.2):
%   - rechts (Seitenzahl):    1,5 cm
%   - links  (Kapiteltitel):  1,5 cm * phi^3 + 3 cm = 9,35 cm
%     Hinweis: ursprueglich 1,5 cm * phi^3 + 2 cm = 8,35 cm (die 2 cm
%     Verlaengerung war ebenfalls eine phi-Groesse, da aus phi^2 = phi + 1
%     die Identitaet 1/phi^2 = 2 - phi folgt, d. h. phi + 1/phi^2 = 2 cm
%     exakt). Die Linie wurde nachtraeglich um weitere 1 cm nach rechts
%     verlaengert, um eine bessere visuelle Ausgewogenheit zwischen
%     Kapiteltitel und Akzent zu erreichen.
% Die Einrueckung der Seitenzahl betraegt 1/phi cm = 0,618 cm.
\newlength{\seitenzahlEinrueckung}\setlength{\seitenzahlEinrueckung}{0.618cm}
\newlength{\akzentlinienBreite}\setlength{\akzentlinienBreite}{1.5cm}
\newlength{\akzentlinienBreiteHaupt}\setlength{\akzentlinienBreiteHaupt}{9.35cm}

% Linker Kopf: Kapiteltitel buendig am linken Seitenrand (keine Einrueckung)
% und linksbuendige Akzentlinie (9,35 cm) darunter.
% Vgl. Leitfaden Kap. 5.2.1: Im Kopfbereich koennen Abschnittsueberschriften
% und Seitennummern erscheinen.
\ihead{%
    \begin{tabular}[b]{@{}l@{}}%
        \headmark \\[-2pt]%
        \textcolor{FOMakzent}{\rule{\akzentlinienBreiteHaupt}{0.8pt}} \\
    \end{tabular}%
}

% Rechter Kopf: Seitenzahl (eingerueckt um 1/phi cm) + rechtsbuendige
% Akzentlinie (1,5 cm) buendig am Seitenrand.
% Kap. 5.2.2: einheitliche Hervorhebungen.
\ohead*{%
    \begin{tabular}[b]{@{}r@{}}%
        \pagemark\hspace*{\seitenzahlEinrueckung} \\[-2pt]%
        \textcolor{FOMakzent}{\rule{\akzentlinienBreite}{0.8pt}} \\
    \end{tabular}%
}
\setkomafont{pagehead}{\normalfont\small}
\setkomafont{pagenumber}{\normalfont\small}

% ----------------------------- Ueberschriften -------------------------------
% Leitfaden Kap. 5.2.3: "Jedes Kapitel, jeder Abschnitt und jeder
% Unterabschnitt erhalten eine Ueberschrift. In der Regel werden sie fett
% und in einem groesseren Schriftgrad als der Haupttext gesetzt. Dabei
% nimmt der Schriftgrad mit zunehmender Gliederungstiefe ab."
\setkomafont{disposition}{\sffamily\bfseries}
\setkomafont{chapter}{\sffamily\bfseries\Large}
\setkomafont{section}{\sffamily\bfseries\large}
\setkomafont{subsection}{\sffamily\bfseries\normalsize}
\setkomafont{subsubsection}{\sffamily\bfseries\normalsize}
% Abstand vor Ueberschriften 12 pt, dahinter 6 pt (Leitfaden Kap. 5.2.2:
% "Der Abstand bei Ueberschriften ... soll vor einem Absatz 12 pt und
% nach einem Absatz 6 pt betragen").
\RedeclareSectionCommand[
    beforeskip=12pt,
    afterskip=6pt
]{chapter}
\RedeclareSectionCommand[
    beforeskip=12pt,
    afterskip=6pt
]{section}
\RedeclareSectionCommand[
    beforeskip=12pt,
    afterskip=6pt
]{subsection}
\RedeclareSectionCommand[
    beforeskip=12pt,
    afterskip=6pt
]{subsubsection}

% Dekadische Gliederung bis Tiefe 4 (Leitfaden Kap. 5.2.3: "Auf
% Nummerierung kann ab einer Gliederungstiefe groesser vier verzichtet
% werden.").
\setcounter{secnumdepth}{4}
\setcounter{tocdepth}{3}

% ----------------------------- Blocksatz & Trennung -------------------------
% Leitfaden Kap. 5.3: "Absaetze werden durchgaengig im Blocksatz gesetzt.
% Dabei ist die Silbentrennung zu verwenden, damit die Wortzwischenraeume
% nicht unverhaeltnismaessig gross werden."
\sloppy
\hyphenpenalty=300
\tolerance=2000
\emergencystretch=10pt

% ----------------------------- Fussnoten ------------------------------------
% Leitfaden Kap. 5.2.2: "Zeilenabstand in Fussnoten ist einzeilig, der
% Schriftgrad 10 pt (Times New Roman) oder 9,5 pt (Arial)." Bei Open Sans
% (sans-serif analog Arial) ist 10 pt eine zulaessige Wahl.
% Vgl. auch Leitfaden Kap. 5.4 zur Verwendung von Fussnoten.
% \deffootnote ist in KOMA-Script direkt verfuegbar (kein zusaetzliches Paket).
\deffootnote[1em]{1.5em}{1em}{\textsuperscript{\thefootnotemark}\ }
\setkomafont{footnote}{\fontsize{10pt}{12pt}\selectfont}
\setkomafont{footnotereference}{\textsuperscript}

% Trennlinie zwischen Fließtext und Fußnoten als Akzentlinie in FOMakzent
% (gleiche Farbe wie die Akzentlinien in Kopfzeile/Titelblatt).
% Laenge 6,354 cm = 1,5 cm · φ³ -- naechste Goldener-Schnitt-Stufe der
% Akzentlinien-Familie nach der 1,5-cm-Linie der Seitenzahl und vor der
% 9,354 cm = 1,5·φ³ + 3 cm langen Kapitellinie. Damit liegt der Wert
% innerhalb des vom Schreibstil geforderten 5--7-cm-Intervalls und ist
% nicht arbiträr, sondern an die uebrigen Linien anschlussfaehig.
% Strichstaerke 0,8 pt (konsistent mit den Header-Linien); Mindestabstand
% von 0,2 cm zwischen Trennlinie und Fußnoten-Fließtext (entspricht der
% kerning-Vorgabe der nachfolgenden Zeile).
\renewcommand*{\footnoterule}{%
    \kern-3pt%
    \hbox{\color{FOMakzent}\vrule width 6.354cm height 0.8pt depth 0pt}%
    \kern 0.2cm%
}

% Durchgaengige Fußnoten-Nummerierung ueber das gesamte Dokument hinweg.
% scrreprt setzt den Fußnoten-Counter standardmaessig zu Beginn jedes
% Kapitels zurueck. Durch das Entkoppeln von "chapter" laeuft der Zaehler
% kontinuierlich von der ersten bis zur letzten Fußnote der Arbeit weiter,
% was Querverweise nachvollziehbar haelt (vgl. Leitfaden Kap. 5.4).
\counterwithout{footnote}{chapter}

% --- Trenner zwischen aufeinanderfolgenden Fußnoten-Markierungen -----------
% Bei aufeinanderfolgenden Fußnoten-Citationen -- etwa
% "...{a}\footcite{b}\footcite{c}" -- kollidieren die Hochziffern im
% Fließtext optisch (Ergebnis "¹²³" statt "¹ ² ³"). Das Paket footmisc
% loest dies mit der Option `multiple`: per kerngesteuerter
% Folge-Erkennung wird zwischen direkt aufeinanderfolgenden
% Hochziffern automatisch ein Trenner gesetzt -- nicht jedoch vor
% Satzzeichen oder vor regulaer mit Leerzeichen abgesetzten Folgewoertern.
% Der Trenner wird ueber \multfootsep konfiguriert; Default ist ein
% Komma, hier auf \space (Leerzeichen) umgestellt.
% Hinweis: Die fuer die biblatex-Kombination noetige Neudefinition von
% \mkbibfootnote (Entfernen des Kern-zerstoerenden \unspace) erfolgt
% weiter unten im Bibliografie-Abschnitt, nachdem biblatex-chicago
% \mkbibfootnote definiert hat.
\usepackage[multiple]{footmisc}
\renewcommand{\multfootsep}{\space}

% ----------------------------- Grafik & Tabellen ----------------------------
% Abbildungen: Leitfaden Kap. 2.3 / 5.7 (Bildunterschrift UNTER der Abbildung,
% fortlaufende Nummerierung, ungerahmt, im Text referenziert).
% Tabellen:    Leitfaden Kap. 2.4 / 5.6 (Titel UEBER der Tabelle, sparsame
% Linienverwendung, fortlaufende Nummerierung).
\usepackage{graphicx}
\graphicspath{{abbildungen/}}
\usepackage{svg}            % Vektor-Einbindung der Modelle/Logo/Unterschrift (benoetigt --shell-escape + Inkscape)
\svgpath{{abbildungen/}}
\svgsetup{inkscapelatex=false}   % Text von Inkscape direkt ins PDF einbetten (verhindert ueberlagerte Text-Artefakte in den BPMN-Modellen)
\usepackage{booktabs}                          % schlanke Tabellenlinien
\usepackage{array}
\usepackage{longtable}
\usepackage{caption}
\usepackage{varwidth}
\captionsetup{
    font={small},
    labelfont={bf,sf},
    labelsep=colon,
    justification=centering,
    skip=6pt                                    % 6 pt Abstand (Kap. 5.2.2)
}
\captionsetup[figure]{position=below}           % Kap. 2.3 / 5.7
\captionsetup[table]{position=above}            % Kap. 2.4 / 5.6

% Einheitliches Makro fuer Quellenangaben unter Abbildungen, das die
% Zweizeiligkeit (Bildunterschrift + Quellenangabe) typografisch stabil und
% projektweit konsistent abbildet (FOM-Leitfaden Kap. 5.7).
% Verwendung in jeder figure-Caption am Ende:
%   \caption[Kurztitel]{Langtitel.\quelleeigen}     % rein eigene Darstellung
%   \caption[Kurztitel]{Langtitel.\quelleeigenki}   % KI-unterstuetzt
%                                                    % erstellte Darstellung
%                                                    % (FOM-Leitfaden Kap. 6.2.4)
%   \caption[Kurztitel]{Langtitel.\quellefremd{Autor (Jahr, S.~X)}}
% Caption-Format: gesamter Beschriftungsblock (Abbildungskennung + Titel +
% Quellenangabe) wird als Box in natuerlicher Breite zentriert unter dem Bild
% gesetzt. Innerhalb der Box ist der Titel zentriert, die Quellenangabe jedoch
% linksbuendig zur Abbildungskennung ausgerichtet (vgl. Anforderung Kap. 5.7).
\DeclareCaptionFormat{quellecentered}{%
  {\centering\begin{varwidth}{\linewidth}\centering #1#2#3\end{varwidth}\par}%
}
\captionsetup[figure]{format=quellecentered}

% Quellenangabe beginnt eine neue, linksbuendige Zeile innerhalb der Box.
\newcommand{\quelleeigen}{\par\vspace{2pt}\raggedright\small Quelle: Eigene Darstellung.}
\newcommand{\quelleeigenki}{\par\vspace{2pt}\raggedright\small Quelle: Eigene Darstellung (KI-generiert).}
\newcommand{\quellefremd}[1]{\par\vspace{2pt}\raggedright\small Quelle: #1.}

% Abbildungs- und Tabellenverzeichnis bündig am linken Seitenrand setzen
% (entspricht der eigenen Darstellung im FOM-Leitfaden, Vorspann S. vi/vii:
% "Abbildung 1: Laufzeitmessung Systemvergleich ........ 22").
\DeclareTOCStyleEntry[indent=0pt,numwidth=3em]{tocline}{figure}
\DeclareTOCStyleEntry[indent=0pt,numwidth=3em]{tocline}{table}

% ----------------------------- Mathematik & Formeln -------------------------
% Leitfaden Kap. 5.8: "Bezeichner von Variablen werden kursiv gesetzt.
% Bezeichner von Konstanten und Operatorsymbole werden aufrecht gesetzt
% ... Eine Gleichung ist eine Aussage und wird mit einem Satzzeichen
% abgeschlossen, sofern der Satzbau dies erfordert. ... Formeln werden
% laufend durchnummeriert."
\usepackage{amsmath,amssymb,amsthm}
\numberwithin{equation}{chapter}                % Gleichungsnummern je Kapitel
% mathastext rendert Math-Buchstaben in der aktuellen Textschrift (Open Sans),
% damit Math und Fliesstext visuell zusammenpassen. Griechische Buchstaben
% bleiben Computer Modern (kein dedizierter Open-Sans-Math-Font verfuegbar).
\usepackage[italic,defaultmathsizes]{mathastext}

% ----------------------------- Quelltexte -----------------------------------
% Leitfaden Kap. 5.9.2 (Quelltexte): "um die Einrueckung zur Geltung zu
% bringen, wird eine Fixtype verwendet ... die Zeilen koennen nummeriert
% werden ... Quelltexte sollten hinreichend durch Kommentare dokumentiert
% sein ... die Quelltexte erhalten Titel".
\usepackage{listings}
\lstset{
    basicstyle=\ttfamily\small,                 % Fixtype (Inconsolata)
    breaklines=true,
    breakatwhitespace=true,
    numbers=left,                                % Zeilennummerierung
    numberstyle=\tiny\color{gray},
    numbersep=8pt,                               % Abstand Nummer-zu-Code
    xleftmargin=2.5em,                           % Listing einruecken: ohne
                                                 % diese Einrueckung wuerden
                                                 % die Zeilennummern in den
                                                 % linken Seitenrand ragen
                                                 % (= negative Position
                                                 % relativ zum Textsatz).
    framesep=4pt,
    frame=single,
    captionpos=t,                                % Titel ueber Listing
    showstringspaces=false,
    tabsize=2
}
\renewcommand{\lstlistingname}{Quelltext}
\renewcommand{\lstlistlistingname}{Quelltextverzeichnis}

% ----------------------------- Listen ---------------------------------------
% Leitfaden Kap. 5.5: Gliederungs-, Beschreibungs- und Aufzaehlungslisten
% mit einheitlicher Symbolverwendung; bei laengeren Listenpunkten
% Blocksatz; geschachtelte Listen erhalten je Stufe ein eigenes Symbol.
\usepackage{enumitem}
\setlist{itemsep=2pt,topsep=4pt}

% ----------------------------- PDF/A & Hyperlinks ---------------------------
% Leitfaden Kap. 5.10.2: "es muss die verwendeten Zeichensaetze eingebettet
% mitbringen ... ein Format gewaehlt werden, das sich moeglichst nah am
% Standard PDF/A orientiert." Das pdfx-Paket erzeugt PDF/A-2b-konforme
% PDFs (Zeichensaetze eingebettet, keine Verschluesselung, kein Schutz
% gegen Drucken/Extraktion). pdfx LAEDT hyperref intern; es darf daher
% kein separates \usepackage{hyperref} davor oder danach erfolgen.
% Erforderlich: Datei main.xmpdata mit Title, Author, Subject, Keywords,
% Publisher, Language.
\usepackage[a-2b]{pdfx}

% Leitfaden Kap. 2.3 / 2.4: "Referenzen sollten mithilfe von Textmarken
% erzeugt werden (bei LATEX: \autoref{label}). ... Die Referenz nennt die
% vollstaendige Bezeichnung: 'Abbildung 3' / 'Tabelle 4'. Referenzen
% werden nicht abgekuerzt." cleveref liefert dies automatisch.
\hypersetup{
    hidelinks,
    bookmarksnumbered=true
}
\usepackage[ngerman]{cleveref}
\crefname{figure}{Abbildung}{Abbildungen}
\crefname{table}{Tabelle}{Tabellen}
\crefname{equation}{Gleichung}{Gleichungen}
\crefname{section}{Abschnitt}{Abschnitte}
\crefname{chapter}{Kapitel}{Kapitel}
\crefname{lstlisting}{Quelltext}{Quelltexte}
\Crefname{lstlisting}{Quelltext}{Quelltexte}

% ----------------------------- Bibliografie (Chicago Notes) -----------------
% Zitierstil: Chicago Manual of Style (17. Aufl.), Notes & Bibliography.
% Belege erscheinen als Fußnoten -- Vollbeleg bei der Erstnennung,
% Kurzbeleg bei Folgenennungen. Das Quellenverzeichnis bleibt
% alphabetisch nach Erstautor sortiert.
% Hinweis zum FOM-Leitfaden: Kap. 6.2 schreibt keinen einzelnen Stil
% verbindlich vor, verlangt aber konsistente Anwendung mit Autor, Jahr und
% Seitenzahl (Kap. 6.2.3). Beide Vorgaben werden durch den Chicago-Notes-
% Stil erfüllt, da jede Fußnote Autor, Werktitel, Jahr und Seitenzahl
% enthält. Kap. 6.3: Quellenverzeichnis alphabetisch nach Autor sortiert.
% Kap. 6.2.4: KI-generierte Inhalte werden separat im KI-Verzeichnis
% nach dem Literaturverzeichnis gefuehrt.
\usepackage[
    notes,                         % Chicago Notes-Bibliography (Fußnoten)
    backend=biber,
    natbib=true,
    sorting=nyt,                   % alphabetisch nach Name/Jahr/Titel
    sortcites=true,
    maxcitenames=2,
    mincitenames=1,
    maxbibnames=99,
    giveninits=true,
    uniquename=init,
    isbn=true,
    doi=true,
    url=true,
    backref=false,
    autocite=footnote              % \autocite -> Fußnote (Chicago-Notes)
]{biblatex-chicago}

% Deutsche Lokalisierung gaengiger Bezeichner (vgl. Leitfaden-Beispiele in
% Kap. 6.2.3); Trennzeichen werden vom Chicago-Stil selbst gesetzt und
% nicht ueberschrieben.
\DefineBibliographyStrings{ngerman}{%
    andothers   = {{et\,al\adddot}},                  % "et al."
    backrefpage = {{vgl\adddotspace S\adddotspace}}
}

% Bibliografie-Datei einbinden
% verzeichnisse/references.bib -- einzige Quelle des Literaturverzeichnisses;
%                   wird ueber Mendeley Reference Manager synchronisiert. Die
%                   globale biblatex-Konfiguration oben (Chicago Notes,
%                   sorting=nyt, ...) gilt dafür.
\addbibresource{verzeichnisse/references.bib}

% --- Mehrfache-Fußnoten-Patch fuer biblatex --------------------------------
% biblatex setzt \mkbibfootnote standardmaessig in der Form
%   \unspace\xspace\bibsentence\footnote{\bibsentence#1\addperiod}
% ein. Das fuehrende \unspace entfernt vorhandene Trailing-Spaces und
% damit auch den Kern-Marker, den footmisc[multiple] hinter jeder
% Fußnoten-Markierung platziert, um aufeinanderfolgende Marken zu
% erkennen. Folge: bei \footcite\footcite-Sequenzen schlaegt die
% Trennererkennung fehl ("¹²" statt "¹ ²").
% Die Neudefinition entfernt das stoerende \unspace und behaelt das
% darauf folgende \bibsentence bei. Da biblatex-chicago Fußnoten-
% Citationen ausschliesslich ueber \mkbibfootnote abwickelt, wirkt
% der Patch zugleich fuer \footcite, \autocite und die Wrapper \zit
% und \zitw.
\renewcommand*{\mkbibfootnote}[1]{%
    \iftoggle{blx@footnote}%
        {#1}%
        {\bibsentence\footnote{\bibsentence#1\addperiod}}%
}

% Filter: Eintraege mit Schluesselwort "ki" -> KI-Verzeichnis (Kap. 6.2.4),
% alle uebrigen -> Hauptliteraturverzeichnis (Kap. 6.3).
\defbibfilter{nichtKI}{not keyword=ki}
\defbibfilter{nurKI}{keyword=ki}

% Hilfsmakros für Fußnoten-Belege (Chicago Notes):
%   \zit[S.~24]{key}  -> Fußnote mit voranstehendem "Vgl. ..."
%   \zitw[S.~24]{key} -> Fußnote ohne "vgl." (wörtliches Zitat)
% Äquivalent zu \footcite (siehe biblatex-chicago); die Kurznamen werden
% aus Lesbarkeitsgründen beibehalten.
\newcommand{\zit}[2][]{\footcite[vgl.][#1]{#2}}
\newcommand{\zitw}[2][]{\footcite[#1]{#2}}

% ----------------------------- Glossar / Abkuerzungen / Symbole ------------
% Leitfaden Kap. 2.5: "Abkuerzungen werden vor der ersten Verwendung formal
% eingefuehrt. Der Begriff wird ausgeschrieben und die abkuerzende
% Schreibweise wird dahinter in Klammern gesetzt." -- glossaries-extra
% automatisiert dies: \gls{api} liefert beim ersten Auftreten "Application
% Programming Interface (API)", danach nur noch "API".
% Leitfaden Kap. 5.1 nennt zusaetzlich "Glossar" und "Verzeichnis von
% Formelzeichen und Indizes" als optionale Verzeichnisse. Beide werden
% hier ueber dieselbe glossaries-extra-Infrastruktur verwaltet:
%   - main           -> Glossar (Fachbegriffe)        (Kap. 5.1)
%   - abbreviations  -> Abkuerzungsverzeichnis        (Kap. 2.5 / 5.1)
%   - symbols        -> Formelzeichenverzeichnis      (Kap. 5.1 / 5.8)
% MUSS NACH hyperref/biblatex geladen werden.
% \makenoidxglossaries (statt \makeglossaries) verzichtet auf externes
% makeglossaries/bib2gls-Tool und erlaubt damit reines pdflatex-Compilieren
% (Overleaf-Default). Sortierung erfolgt beim Drucken via TeX.
% Wichtig: glossaries-extra laedt die Style-Files (altlist, longtable, ...)
% nicht automatisch. Ohne `stylemods=...` wuerde \printnoidxglossary[
% style=altlist] mit "Undefined control sequence" bzw. "Illegal parameter
% number in definition of \csname\endcsname" abbrechen. Der custom Style
% `fomabbr` (Abkuerzungsverzeichnis) ist davon nicht betroffen, weil er
% direkt per \newglossarystyle definiert ist.
\usepackage[%
    abbreviations,%             Abkuerzungs-Liste aktiv
    symbols,%                   Formelzeichen-Liste aktiv (Kap. 5.1 / 5.8)
    nonumberlist,%              keine Seitenzahlen im Verzeichnis
    nopostdot,%                 kein Punkt am Eintrags-Ende
    toc,%                       Eintrag im Inhaltsverzeichnis
    stylemods=list%             altlist/list-Styles fuer das Glossar laden
]{glossaries-extra}
\setabbreviationstyle[abbreviation]{long-short}
\makenoidxglossaries

% Custom-Style "fomabbr" fuer das Abkuerzungsverzeichnis:
%   - longtable startet bei 0pt (flush-left zum Seitenrand)
%   - linke Spalte 2,5 cm breit, Eintraege fett gesetzt
%   - 0,7 cm Spaltenabstand
%   - rechte Spalte raggedright, restliche Textbreite
% Dadurch werden Eintraege wie "API" / "REST" linksbuendig zum Seitenrand
% gesetzt; die Erlaeuterung beginnt einheitlich 0,7 cm hinter der
% Kurzform-Spalte.
\newglossarystyle{fomabbr}{%
    \renewenvironment{theglossary}{%
        \setlength{\LTleft}{0pt}%
        \setlength{\LTright}{\fill}%
        \begin{longtable}{%
            @{}>{\bfseries}p{2.5cm}@{\hspace{0.7cm}}%
            >{\raggedright\arraybackslash}p{\dimexpr\textwidth-3.2cm\relax}@{}%
        }%
    }{%
        \end{longtable}%
    }%
    \renewcommand*{\glossaryheader}{}%
    \renewcommand*{\glsgroupheading}[1]{}%
    \renewcommand*{\glossentry}[2]{%
        \glsentryitem{##1}\glstarget{##1}{\glossentryname{##1}}%
        & \glossentrydesc{##1}\glspostdescription \tabularnewline
    }%
    \renewcommand*{\subglossentry}[3]{%
        & \glossentrydesc{##2}\glspostdescription \tabularnewline
    }%
    \renewcommand*{\glsgroupskip}{}%
}

% Custom-Style "fomgloss" fuer das Glossar (Fachbegriffe). Analog zu
% `fomabbr` ohne Abhaengigkeit von glossary-list/longtable -- damit ist
% das Glossar auch dann darstellbar, wenn `stylemods` keine Wirkung
% entfaltet. Aufbau:
%   Begriff (fett, eigener Absatz)
%   Beschreibung (eingerueckt, Blocksatz)
\newglossarystyle{fomgloss}{%
    \renewenvironment{theglossary}{%
        \setlength{\parindent}{0pt}%
        \setlength{\parskip}{6pt plus 1pt}%
    }{}%
    \renewcommand*{\glossaryheader}{}%
    \renewcommand*{\glsgroupheading}[1]{}%
    \renewcommand*{\glossentry}[2]{%
        \glsentryitem{##1}\glstarget{##1}{%
            \textbf{\glossentryname{##1}}}\par
        \nopagebreak\noindent\hangindent=1.5em\hangafter=0
        \glossentrydesc{##1}\glspostdescription\par
    }%
    \renewcommand*{\subglossentry}[3]{%
        \glsentryitem{##2}\glstarget{##2}{%
            \textbf{\glossentryname{##2}}}\par
        \nopagebreak\noindent\hangindent=1.5em\hangafter=0
        \glossentrydesc{##2}\glspostdescription\par
    }%
    \renewcommand*{\glsgroupskip}{}%
}

\loadglsentries{verzeichnisse/abkuerzungen-defs}
\loadglsentries{verzeichnisse/glossar-defs}
\loadglsentries{verzeichnisse/formelverzeichnis-defs}

% ----------------------------- Sonstiges ------------------------------------
% Farben: Leitfaden Kap. 5.2.2 verlangt Hervorhebungen "sparsam und
% einheitlich". Die Akzentfarbe ist deshalb auf Titelblatt + Header-
% Akzentlinie beschraenkt; Fliesstext bleibt schwarz/weiss.
\usepackage{xcolor}
\definecolor{FOMakzent}{HTML}{239F91}        % FOM-Petrolgruen (Logo)
\definecolor{FOMakzentDunkel}{HTML}{1B7A6F}  % dunklere Variante

% Lokalisierung: Komma als Dezimaltrenner ist im deutschen Schriftsatz
% Standard (vgl. Leitfaden Kap. 5.6 Tabelle 1, "Wechselkurs ... 1,13").
\usepackage{eurosym}
\usepackage{siunitx}
\sisetup{
    locale = DE,
    output-decimal-marker = {,},
    group-separator = {\,},
    detect-weight = true,
    detect-family = true
}

% Leitfaden Kap. 7.4: "Hurenkinder und Schusterjungen (... einzelne Zeilen
% eines Absatzes am Anfang oder am Ende einer Seite)" sollen vermieden
% werden. Hohe Penalties unterdruecken solche Brueche.
\clubpenalty=10000
\widowpenalty=10000
\displaywidowpenalty=10000

% ============================================================================
%                              METADATEN
% Pflichtangaben des Titelblatts (Leitfaden Kap. 5.1, Anhang I):
% Identifikation (Name, Matrikelnummer), Titel, Art der Arbeit, Hochschule
% und Studiengang, Erstgutachter, Datum/Ort der Fertigstellung.
% ============================================================================
\newcommand{\arbeitTitel}{OpenClaw als Trading Agent}
\newcommand{\arbeitUntertitel}{AI Agent zur Nutzung einer SMA Strategie im Derivate-Handel}
\newcommand{\arbeitArt}{Projektarbeit} % Bachelor-Thesis oder Master-Thesis / Seminararbeit
\newcommand{\arbeitGrad}{Bachelor of Science (B.Sc.)}
\newcommand{\arbeitStudiengang}{Informatik}
\newcommand{\arbeitHochschulzentrum}{Digitales Live-Studium (DLS)}
% Autoren alphabetisch nach Nachname (Gruppenarbeit)
\newcommand{\autorVorname}{Alexander}
\newcommand{\autorNachname}{Fischer}
\newcommand{\autorMatrikel}{825091}
\newcommand{\autorZweiVorname}{Rolland}
\newcommand{\autorZweiNachname}{Kiss}
\newcommand{\autorZweiMatrikel}{799553}
\newcommand{\erstgutachter}{Prof. Dr. Klemens Waldhör}
\newcommand{\zweitgutachter}{-} % optional
% Vor der finalen Abgabe das tatsaechliche Datum festsetzen, damit das
% Datum nicht bei spaeteren Re-Compiles unbeabsichtigt mitwandert (z. B.
% \newcommand{\abgabedatum}{15.\,Juni 2026}).
\newcommand{\abgabedatum}{31.\,Juli 2026}
\newcommand{\arbeitOrt}{Köln}        % Ort fuer Unterschriften (Kap. 7.2)

% PDF-Metadaten: Title/Author/Subject/Keywords werden bei PDF/A primaer
% aus main.xmpdata gelesen (vom pdfx-Paket). Die folgenden Werte spiegeln
% sie fuer Nicht-PDF/A-Builds wider.
\hypersetup{
    pdftitle={\arbeitTitel},
    pdfauthor={\autorVorname{} \autorNachname{}, \autorZweiVorname{} \autorZweiNachname},
    pdfsubject={\arbeitArt{} - FOM Hochschule},
    pdfkeywords={FOM, \arbeitStudiengang, \arbeitArt}
}

% ============================================================================
%                              DOKUMENT
% ============================================================================
\begin{document}

% --- Vorspann (roemische Seitenzaehlung) ------------------------------------
% Leitfaden Kap. 5.2.1: "Die Seiten des Vorspanns werden im Allgemeinen
% roemisch nummeriert. Danach wird auf arabische Seitennummerierung
% beginnend mit Eins umgestellt." Im Beispiel des Leitfadens werden
% kleinroemische Ziffern (i, ii, iii) verwendet.
% Titelblatt = Seite i, Inhaltsverzeichnis = Seite ii (Sperrvermerk wird
% gemaess Leitfaden Kap. 7.2 NICHT mitgezaehlt).
\pagenumbering{roman}
\pagestyle{plain}

%  =============================================================================
%  -> Abschnitt "titelblatt/titelblatt.tex"
%  =============================================================================

\begin{titlepage}
\thispagestyle{empty}

% --- Kopfzeile: Logo links, Hochschultext rechts ---------------------------
\noindent
\begin{minipage}[c]{0.27\textwidth}
    \includesvg[height=2.2cm]{fom-logo}
\end{minipage}\hfill
\begin{minipage}[c]{0.7\textwidth}
    \raggedleft
    {\large\sffamily FOM Hochschule für Oekonomie \& Management}\\[0.15cm]
    {\large\sffamily Hochschulzentrum \arbeitHochschulzentrum}
\end{minipage}

\vspace{0.4cm}
{\color{FOMakzent}\rule{\textwidth}{0.8pt}}

\begin{center}
\vspace{1cm}

% --- Art der Arbeit ---------------------------------------------------------
{\Large\sffamily\bfseries \arbeitArt}\\[0.25cm]
{\large\sffamily im Studiengang \arbeitStudiengang}\\[0.9cm]

zur Erlangung des Grades eines\\[0.25cm]
{\Large\sffamily\bfseries \arbeitGrad}\\[0.9cm]

über das Thema\\[0.5cm]

% --- Titel der Arbeit -------------------------------------------------------
{\huge\sffamily\bfseries \arbeitTitel}\\[0.3cm]
\ifx\arbeitUntertitel\empty\else
    {\Large\sffamily \arbeitUntertitel}\\[0.8cm]
\fi
\vspace{0.7cm}

von\\[0.3cm]
{\LARGE\sffamily \autorVorname{} \autorNachname}\\[1cm]

\vfill

% --- Akzentlinie + Pruefungsangaben ----------------------------------------
{\color{FOMakzent}\rule{\textwidth}{0.8pt}}\\[0.4cm]
\begin{flushleft}
\renewcommand{\arraystretch}{1.35}
\begin{tabular}{@{}p{4.5cm}p{9cm}@{}}
Erstgutachter      & \erstgutachter \\
\ifx\zweitgutachter\empty\else
Zweitgutachter     & \zweitgutachter \\
\fi
Matrikelnummer     & \autorMatrikel \\
Abgabedatum        & \abgabedatum \\
\end{tabular}
\end{flushleft}

\end{center}
\end{titlepage}
\cleardoublepage

%  =============================================================================
%  Ende
%  =============================================================================


% Sperrvermerk (Leitfaden Kap. 7.2). Nur bei vertraulichen Arbeiten aktivieren;
% die Datei sperrvermerk/sperrvermerk.tex ist im Projekt derzeit nicht
% vorhanden und muesste zuvor angelegt werden:
%\input{sperrvermerk/sperrvermerk}

% Inhaltsverzeichnis
\pagestyle{scrheadings}
\tableofcontents
\cleardoublepage

% Verzeichnisse (Leitfaden Kap. 5.1: "Die aufgefuehrten Verzeichnisse
% erhalten keine Gliederungsnummern, werden aber im Inhaltsverzeichnis
% aufgefuehrt."). Reihenfolge alphabetisch nach Verzeichnis-Bezeichnung:
%   Abbildungs- -> Abkuerzungs- -> Glossar.
% Hinweis: KEIN explizites \clearpage zwischen den Verzeichnissen, da
% jedes Verzeichnis intern \chapter*{...} aufruft und damit selbst eine
% neue Seite startet. Ein zusaetzliches \clearpage erzeugt eine leere
% Zwischenseite mit der Kopfmarke des Vorgaenger-Verzeichnisses.
% Hinweis zu entfernten/ausgeblendeten Verzeichnissen:
%   - Algorithmenverzeichnis: vollstaendig entfernt.
%   - Mathematisches Formelverzeichnis: vollstaendig entfernt.
%   - Symbolverzeichnis: vollstaendig entfernt (Quelldatei nicht vorhanden).
%   - Tabellenverzeichnis: ausgeblendet (nur auskommentiert, s. u.).
% Alle Verzeichnisse liegen einheitlich unter verzeichnisse/ und werden
% gleichartig per \input eingebunden (das Abbildungsverzeichnis kapselt
% \listoffigures in verzeichnisse/abbildungsverzeichnis.tex).
% =============================================================================
%  Datei-Hinweise & Konfiguration: siehe config.md
%  -> Abschnitt "verzeichnisse/abbildungsverzeichnis.tex"
%
%  Abbildungsverzeichnis (Leitfaden Kap. 2.3 / 5.1 / 5.7). Ausgelagert in eine
%  eigene Verzeichnis-Datei, damit alle Verzeichnisse des Vorspanns einheitlich
%  unter verzeichnisse/ liegen und in main.tex gleichartig per \input
%  eingebunden werden. \listoffigures erzeugt das Verzeichnis aus den
%  \caption-Kurztiteln; der TOC-Eintrag folgt aus der KOMA-Option listof=totoc.
% =============================================================================

\listoffigures

%  Kopfzeilen-Marke explizit setzen (einseitiges Verzeichnis; gleiche
%  Konvention wie in den uebrigen Verzeichnis-Dateien).
\markboth{Abbildungsverzeichnis}{Abbildungsverzeichnis}

%  =============================================================================
%  Ende
%  =============================================================================
         % Kap. 2.3 / 5.1 / 5.7
% =============================================================================
%  Datei-Hinweise & Konfiguration: siehe config.md
%  -> Abschnitt "verzeichnisse/abkuerzungen.tex"
% =============================================================================

\glsaddall[types=abbreviations]
\printnoidxglossary[%
    type=abbreviations,%
    title={Abkürzungsverzeichnis},%
    style=fomabbr%
]

%  Kopfzeilen-Marke unter \manualmark explizit setzen (einseitiges Verzeichnis;
%  Kapitelwechsel-Logik siehe main.tex / config.md).
\markboth{Abkürzungsverzeichnis}{Abkürzungsverzeichnis}

%  =============================================================================
%  Ende
%  =============================================================================
                  % Kap. 2.5 / 5.1
%  =============================================================================
%  Datei-Hinweise & Konfiguration: siehe config.md
%  -> Abschnitt "verzeichnisse/formelzeichen.tex"
%     (Definitionen: verzeichnisse/formelverzeichnis-defs.tex)
%
%  Formelzeichenverzeichnis (Leitfaden Kap. 5.1 / 5.8): alle in den
%  nummerierten Gleichungen des Hauptteils verwendeten Formelzeichen mit
%  Kurzerklaerung. \glsaddall nimmt saemtliche Symbol-Eintraege auf
%  (unabhaengig von \gls-Verwendung im Fliesstext, da Formelzeichen in
%  Gleichungen als reine Mathematik gesetzt werden).
%  =============================================================================

\glsaddall[types={symbols}]
\printnoidxglossary[
    type=symbols,
    style=fomabbr,
    title={Formelzeichenverzeichnis}
]

%  =============================================================================
%  Ende
%  =============================================================================
                 % Kap. 5.1 / 5.8
% =============================================================================
%  Datei-Hinweise & Konfiguration: siehe config.md
%  -> Abschnitt "verzeichnisse/glossar.tex"
% =============================================================================

\glsaddall[types={main}]
\printnoidxglossary[%
    type=main,%
    title={Glossar},%
    style=fomgloss,%
    sort=word%
]

%  Kopfzeilen-Marke unter \manualmark explizit setzen (einseitiges Verzeichnis;
%  Kapitelwechsel-Logik siehe main.tex / config.md).
\markboth{Glossar}{Glossar}

%  =============================================================================
%  Ende
%  =============================================================================
                       % Kap. 5.1
%\listoftables                                      % Tabellenverzeichnis ausgeblendet
\clearpage

% Quelltextverzeichnis bewusst nicht aktiviert. Begruendung:
%   - Leitfaden Kap. 5.1 stuft das Verzeichnis als optional ein
%     ("kann sinnvoll sein, wenn ...").
%   - Saemtliche Quelltexte sind in einem einzigen Anhangskapitel gebuendelt
%     und durch dessen Abschnittsstruktur (I.I, I.II, I.III) bereits navigierbar.
%   - Die zusaetzliche Verzeichnisseite waere in dieser Konstellation eine
%     redundante Wiederholung des Anhang-TOC.
%   - Querverweise via \cref{lst:...} ermoeglichen die direkte Navigation.
% Bei Bedarf — etwa wenn Listings ueber den Hauptteil verteilt eingesetzt
% werden — die folgenden Zeilen aktivieren (KOMA fuegt mit listof=totoc
% automatisch einen TOC-Eintrag hinzu, daher KEIN \addcontentsline noetig):
%\lstlistoflistings
%\cleardoublepage

% --- Hauptteil (arabische Seitenzaehlung ab Kapitel 1) ----------------------
% Leitfaden Kap. 5.2.1: arabische Seitennummerierung beginnt mit "1" am
% ersten Kapitel des Hauptteils.
% Aufbau gemaess Kap. 5.1: Einleitung -> Grundlagen -> Umsetzung -> Schluss.
\pagenumbering{arabic}
\setcounter{page}{1}

%  =============================================================================
%  -> Abschnitt "kapitel/01_einleitung.tex"
%  =============================================================================

\chapter{Einleitung}
\label{cha:einleitung}

\section{Problemstellung und Motivation}
\label{sec:problemstellung}

xxx

\section{Zielsetzung und Relevanz der Arbeit}
\label{sec:zielsetzung}

xxx

\newpage
\section{Vorgehen und Methodik}
\label{sec:methodik}

xxx

%  =============================================================================
%  Ende
%  =============================================================================

%  =============================================================================
%  -> Abschnitt "kapitel/02_grundlagen.tex"
%  =============================================================================

\chapter{Theoretische Grundlagen zu KI-Agenten}
\label{cha:grundlagen}

xxx

\section{Begriffe und Ziele}
\label{sec:grundlagen-ziele}

xxx

\section{Einordnung}
\label{sec:grundlagen-einordnung}

xxx

\section{Abgrenzung}
\label{sec:grundlagen-abgrenzung}

xxx

%  =============================================================================
%  Ende
%  =============================================================================

%  =============================================================================
%  -> Abschnitt "kapitel/03_umsetzung.tex"
%
%  Gliederung an der KI-Fallstudie ausgerichtet:
%  Architektur -> Strategie -> Implementierung -> Evaluation -> Wuerdigung.
%  =============================================================================

\chapter{Fallstudie: OpenClaw als Trading Agent}
\label{cha:umsetzung}

Nachdem das vorangegangene Kapitel mit dem Periodensystem der Künstlichen Intelligenz
einen Ordnungsrahmen bereitgestellt hat, überträgt das folgende Kapitel diesen Rahmen
auf den konkreten Gegenstand der Arbeit. Der untersuchte Agent verfolgt eine auf
gleitenden Durchschnitten beruhende Handelsstrategie und löst die daraus abgeleiteten
Entscheidungen eigenständig als Ordervorgänge über eine Bankschnittstelle aus. Damit
schließt er den eingangs beschriebenen Kreis aus Wahrnehmung, Schlussfolgerung und
Handlung vollständig und lässt sich als Verbindung mehrerer KI-Elemente im Sinne des
Periodensystems auffassen.

Auf der Ebene der Wahrnehmung stützt sich der Agent auf die kontinuierliche Erfassung
und Aufbereitung von Marktdaten, deren Qualität über eine unabhängige Zweitquelle
gegengeprüft wird, sowie auf die Erkennung von Trendmustern in der Kurszeitreihe. Auf
der Ebene der Schlussfolgerung tritt ein Gremium aus drei großen Sprachmodellen als
bewertende Instanz hinzu, das die regelbasiert vorbereitete Entscheidung jeweils
unabhängig prüft, mit einer nachvollziehbaren Begründung versieht und dessen Voten nach
einer strikten Mehrheitsregel aggregiert werden; deterministische Vorprüfungen
entscheiden harte Ausschlusskriterien bereits vor dem ersten Modellaufruf. Die
Modellinferenz erfolgt über eine einheitliche Schnittstelle, die den Wechsel zwischen
cloud-gehosteter Ausführung und dem lokalen Betrieb auf eigener Hardware als reine
Konfigurationsentscheidung belässt. Auf der Ebene der Reaktion schließlich ordnet
sich das Vorhaben unmittelbar dem Element der prozesssteuernden Kontrolle zu, das der
zugrunde gelegte Rahmen ausdrücklich mit dem Börsenhandel als Anwendung
verbindet.\zit[S. 18]{waldhoerKI2025} Ergänzend sorgt eine Kommunikationskomponente für
die Rückbindung an den Menschen: Sie übermittelt die tägliche Freigabe der Banksitzung,
berichtet jeden Entscheidungslauf und stellt einen jederzeit wirksamen Not-Aus bereit;
ein öffentliches, redigiertes Dashboard macht die Entscheidungs- und Wertentwicklung
zusätzlich transparent.

Kennzeichnend für die gewählte Gestaltung ist weniger die Breite der eingesetzten
Fähigkeiten als vielmehr deren bewusste Begrenzung. Die Sprachmodelle werden nicht zur
Prognose künftiger Kurse und nicht zur erfahrungsbasierten Selbstveränderung eingesetzt;
die entsprechenden Elemente des Periodensystems bleiben bewusst unbesetzt, um die
Nachvollziehbarkeit und Reproduzierbarkeit des Systemverhaltens zu wahren. Das Gremium
kann Handelsvorschläge ausschließlich verhindern, verkleinern oder zeitlich einordnen,
niemals erzeugen oder in ihrer Richtung umkehren. Flankiert wird diese begrenzte
Diskretion durch einen mehrstufig abgesicherten Orderpfad: Ein zweiphasiges
Ausführungsjournal verhindert die Doppelausführung desselben Entscheids, ein
Orderbuchabgleich klärt unklare Ausführungszustände, die Scharfschaltung auf reale
Orders erfolgt ausschließlich kontospezifisch, und jeder unklare Zustand endet in einem
alarmierten Handelsstopp statt in einem automatischen Wiederholungsversuch. Die
eigentliche gestalterische Aussage der Arbeit liegt somit in der sicherheitsorientierten
Einschränkung des KI-Einsatzes, deren Ausprägung im weiteren Verlauf des Kapitels
dargelegt wird.

Der Aufbau des Kapitels folgt dem Weg von der Struktur zur Bewertung. Zunächst wird die
Systemarchitektur des Agenten mit ihren Komponenten und Datenflüssen beschrieben. Darauf
aufbauend wird die Handelsstrategie formal gefasst und ihre Umsetzung in Quelltext
dargestellt. Die anschließende Evaluation ordnet die Ergebnisse anhand geeigneter
Kennzahlen ein und benennt die methodischen Grenzen der Untersuchung. Eine kritische
Würdigung der mit dem autonomen Handel verbundenen Risiken und der vorgesehenen
Absicherungsmechanismen beschließt das Kapitel.

\section{Systemarchitektur des Agenten}
\label{sec:umsetzung-architektur}

Die Architektur des Agenten folgt einem durchgängigen Gestaltungsprinzip: Module
kommunizieren nicht durch direkte Aufrufe, sondern ausschließlich über gemeinsamen,
persistierten Zustand in Form von SQLite-Datenbanken und synchronisierten Dateien.
Diese lose Kopplung erlaubt es, jede Komponente isoliert zu testen, auszutauschen und
im Fehlerfall neu zu starten, ohne dass sich Störungen über Prozessgrenzen hinweg
fortpflanzen. \Cref{fig:zielarchitektur} zeigt die Zielarchitektur mit den bereits
umgesetzten und den geplanten Bestandteilen.

% Grafik: abbildungen/zielarchitektur.svg — Einbindung als Vektorgrafik via
% \includesvg (svg-Paket, \svgpath{{abbildungen/}}, analog fom-logo).
% Das SVG ist attributbasiert gestylt (kein <style>-Block), damit die
% Inkscape-Konvertierung von Overleaf verlustfrei arbeitet.
\begin{figure}[htbp]
  \centering
  \includesvg[width=\textwidth]{zielarchitektur}
  \caption[Zielarchitektur des autonomen Handelsagenten]{Zielarchitektur des
    autonomen Handelsagenten. Durchgezogene Elemente sind umgesetzt und in Betrieb,
    gestrichelte Elemente geplant.\quelleeigenki}
  \label{fig:zielarchitektur}
\end{figure}

Das produktive System verteilt sich auf zwei virtuelle Maschinen der Google Cloud
Platform. Die \emph{collector-vm} sammelt rund um die Uhr Marktdaten des S\&P~500:
Rohticks werden über einen WebSocket-Datenstrom sowie ein ergänzendes Polling für
Futures-Randzeiten erfasst, durch einen Sanitätsfilter geprüft und zu Minuten- und
Tagesbars verdichtet. Der jeweils jüngste Tagesschluss wird gegen die offizielle
Zweitquelle der Federal Reserve Bank of St.\ Louis (FRED) abgeglichen; Divergenzen
oberhalb einer Schwelle lösen einen Alarm aus. Die \emph{agent-vm} erhält die
Datenbasis als minütlich synchronisierte, atomar ersetzte Replik und beherbergt den
eigentlichen Entscheidungs- und Orderpfad: die deterministische Strategieauswertung,
den mehrstufigen Sprachmodell-Audit, die Orchestrierung der Orderausführung sowie die
Anbindung an die Bankschnittstelle. Beide Maschinen sind wechselseitig überwacht;
ein Totmannschalter auf der collector-vm alarmiert den Betreiber über den
Messaging-Kanal, sobald der Lebenszeichen-Stempel der agent-vm veraltet.

Für den Dauerbetrieb ist die Datensammlung gezielt gehärtet. Bricht der
WebSocket-Datenstrom still ab, erkennt ein Wächter die ausbleibenden Ticks und
erzwingt eine Neuverbindung mit ansteigenden Wartezeiten; der Polling-Prozess für
die Futures-Randzeiten startet sich nach Ausnahmen selbst neu, und ein geordneter
Shutdown stellt sicher, dass keine halb geschriebenen Datensätze zurückbleiben.
Ein Plattenplatz-Wächter, eine wöchentliche Integritätsprüfung der Datenbank und
eine kontinuierliche Replikation des Datenbestands in einen Cloud-Speicher ergänzen
die Vorsorge; Alarme beider Maschinen laufen über eine gemeinsame Brücke im
Messaging-Kanal zusammen. Diese Betriebsdetails sind für die Fragestellung der
Arbeit nicht nebensächlich: Ein Agent, der täglich unbeaufsichtigt Geld bewegen
soll, ist nur so vertrauenswürdig wie die Infrastruktur, die seine Datenbasis und
seine Erreichbarkeit garantiert.

Die Sprachmodellinferenz ist bewusst hinter einer einheitlichen Schnittstelle
(Ollama) gekapselt. Im gegenwärtigen Betrieb werden drei cloud-gehostete Modelle
unterschiedlicher Anbieter befragt; die Migration auf lokal betriebene Modellvarianten
eines Mac~mini ist als reiner Konfigurationswechsel angelegt, da Modellbezeichner,
Endpunkt und Zeitbudgets in der zentralen Konfigurationsdatei liegen und kein
Quelltext angepasst werden muss. Damit bleibt das in \cref{sec:zielsetzung}
formulierte Ziel der lokalen, weitgehend kostenfreien Betreibbarkeit strukturell
gesichert, während der laufende Betrieb von der Verfügbarkeit leistungsfähiger
gehosteter Modelle profitiert.

Eine dritte, strikt entkoppelte Zone bildet das öffentliche Dashboard. Die
collector-vm exportiert minütlich redigierte JSON-Feeds des Entscheidungsjournals
-- Signale, Voten, Orders und Depotwertverlauf, jedoch keinerlei Konto- oder
Zugangsdaten -- in einen privaten Cloud-Storage-Bucket; ein zustandsloser
Cloud-Run-Dienst rendert daraus die Weboberfläche. Die handelnden Maschinen selbst
bleiben ohne offenen eingehenden Port. Im Ordnungsrahmen des Periodensystems der
Künstlichen Intelligenz verbindet die Architektur damit Elemente der Wahrnehmung
(Daten- und Mustererfassung), der Schlussfolgerung (Entscheiden und Erklären) und
der Reaktion, deren Zentrum das Element der prozesssteuernden Kontrolle
bildet.\zit[S.~18]{waldhoerKI2025}

\section{SMA-Handelsstrategie}
\label{sec:umsetzung-strategie}

Die Handelsentscheidungen des Agenten entstammen einer deterministischen
Regimestrategie auf Basis gleitender Durchschnitte. Für die Tagesschlusskurse
$P(t)$ des S\&P~500 ist der einfache gleitende Durchschnitt definiert als
\begin{equation}
  \mathrm{SMA}_n(t) \;=\; \frac{1}{n}\sum_{i=0}^{n-1} P(t-i),
  \qquad n \in \{50,\, 200\}.
  \label{eq:sma}
\end{equation}
Das Handelsregime ergibt sich aus dem Verhältnis der schnellen zur langsamen
Durchschnittslinie:
\begin{equation}
  R(t) \;=\;
  \begin{cases}
    \text{LONG},  & \text{falls } \mathrm{SMA}_{50}(t) > \mathrm{SMA}_{200}(t),\\
    \text{SHORT}, & \text{sonst.}
  \end{cases}
  \label{eq:regime}
\end{equation}
Die Entscheidung für den Handelstag $t{+}1$ nutzt ausschließlich Schlusskurse bis
einschließlich Tag $t$, sodass kein Blick in die Zukunft (Look-Ahead) möglich ist.
Das Regime wird über eine zentrale Zuordnungsdatei auf konkrete Handelsprodukte
abgebildet -- ein Faktor-Zertifikat mit zweifachem Hebel in Long- beziehungsweise
Short-Richtung --, wobei ein Katalog mit Eignungsregeln (Mindestkurs, maximaler
Spread, Hebelobergrenze) die Produktauswahl begrenzt.

Der Kapitaleinsatz je Handelstag wird volatilitätsabhängig skaliert:
\begin{equation}
  S(t) \;=\; \operatorname{clamp}\!\Bigl(K \cdot \min\bigl(1,\,
  \sigma^{*}/\sigma_{20}(t)\bigr),\; S_{\min},\; K\Bigr),
  \label{eq:sizing}
\end{equation}
mit dem Basiskapital $K = 2.000\,€$, der realisierten
20-Tage-Volatilität $\sigma_{20}$ der Index-Tagesrenditen, der Zielvolatilität
$\sigma^{*} = 1{,}0\,\%$ pro Tag und der Mindestgröße $S_{\min} = 1.000\,€$. Die
Obergrenze $K$ ist als harter, weder über die Konfiguration noch durch die
Sprachmodelle veränderbarer Deckel implementiert.

Die Ausführungsregeln sind ebenfalls deterministisch festgelegt. Gehandelt wird im
Fenster der US-Kassasitzung mit drei diskreten Einstiegsoptionen, unter denen das
Prüfgremium wählen darf: dem Einstieg kurz nach der Eröffnungsauktion, einem
bestätigungsbasierten Einstieg nach den ersten fünfzehn Handelsminuten oder einem
bewusst verzögerten Einstieg, der die Eröffnungsvolatilität meidet. Der Ausstieg
folgt dem zuerst erreichten Kriterium aus Verluststopp (drei Prozent), Gewinnziel
(zwei Prozent) und Fensterschluss. Es gilt ein hartes Tagesbudget von höchstens
einem Kauf und einem Verkauf, und sämtliche Orders werden ausschließlich limitiert
aufgegeben, da die Datenfrische im verteilten Aufbau eine Slippage-Kontrolle über
den Limitpreis erfordert.

Die ökonomische Hürde dieser Taktung lässt sich vorab beziffern: Ein Round-Trip
kostet im Direkthandel 7{,}80\,€ an Fixgebühren, hinzu kommt der Geld-Brief-Spread
des Zertifikats. Bei einem Einsatz von 2.000\,€ muss sich der Index folglich um
etwa 0{,}3\,\% in Regimerichtung bewegen, bevor ein Handelstag den Break-even
erreicht -- eine erreichbare, aber keineswegs sichere Schwelle, deren Konsequenzen
die Evaluation quantifiziert. Die Mindestgröße von 1.000\,€ folgt umgekehrt aus der
Gebührenstruktur: Darunter dominiert die Fixkostenlast jede realistische
Tagesbewegung.

Als extern validierte Alternative ist die Regimestrategie \emph{smatrend}
implementiert, die statt eines Kreuzungssignals den Indexstand gegen einen einzelnen
langfristigen Durchschnitt stellt:
\begin{equation}
  R'(t) \;=\;
  \begin{cases}
    \text{LONG}, & \text{falls } P(t) > \mathrm{SMA}_{325}(t),\\
    \text{CASH}, & \text{sonst,}
  \end{cases}
  \label{eq:smatrend}
\end{equation}
und nur bei einem Regimewechsel handelt, ansonsten die Position hält. Beide
Strategien folgen demselben Plugin-Kontrakt (Signalserie, deklarierter Datenbedarf,
Mindesthistorie, Entscheidungsintervall), sodass der geplante A/B-Vergleich auf
getrennten Depots ohne Eingriff in Auditor, Validator oder Broker möglich ist.

\section{Implementierung}
\label{sec:umsetzung-implementierung}

Grundlage jeder Entscheidung ist die Datenschicht, die einem strikten Prinzip folgt:
Gespeichert werden Rohticks, aus denen Bars abgeleitet werden -- nie umgekehrt.
Damit bleibt jede Verdichtung reproduzierbar, und Auffälligkeiten lassen sich bis
zum einzelnen Tick zurückverfolgen. Ein Sanitätsfilter verwirft unplausible Werte
nicht, sondern markiert sie, sodass die Provenienz erhalten bleibt. Lücken im
Datenstrom werden nach ihrer Ursache klassifiziert; ein empirischer Befund des
Betriebs unterstreicht, warum das nötig ist: Die reguläre Stille des Index
außerhalb der US-Kassasitzung hätte an jedem Montagmorgen fälschlich als
Datenfehler gegolten und den Handel strukturell blockiert. Erst die Unterscheidung
zwischen legitimer Sitzungspause und echtem Verbindungsverlust machte das
Qualitätskriterium betriebstauglich -- ein Beispiel dafür, dass sich die
Schwellwerte eines autonomen Systems nicht am Schreibtisch, sondern nur gegen
reale Betriebsdaten kalibrieren lassen.

Der tägliche Zyklus beginnt mit der Erzeugung eines strukturierten
Entscheidungsdokuments (\emph{DecisionRequest}): Die Strategieauswertung liefert
Regime, Indikatorwerte, Signalalter, Volatilitätskennzahlen, Datenqualitätsmetriken
und einen deterministisch berechneten Kostenkontext (Round-Trip-Fixkosten und deren
prozentuale Last je Größenstufe). Ergänzend setzt sie Prüfmarken
(\emph{Audit-Flags}), von denen ein Teil als harte Ausschlusskriterien definiert ist
-- etwa veraltete Daten, jüngste Verbindungsverluste des Datenstroms oder ein
fehlendes Handelsinstrument. Ein vorgeschaltetes Pre-Gate wertet diese harten
Kriterien vor jedem Modellaufruf aus: Steht das Ergebnis deterministisch fest, wird
die Ablehnung ohne Sprachmodellbeteiligung journalisiert; die Modelle werden nur
befragt, wenn ihre Bewertung das Ergebnis noch verändern kann.

Im Auditschritt erhält jedes der drei konfigurierten Sprachmodelle denselben
DecisionRequest mit einem Systemauftrag als kritischer Risikoprüfer. Die Antwort ist
über ein erzwungenes JSON-Schema strukturiert und wird bei Temperatur null erzeugt;
jedes Modell muss neben der Zustimmung oder Ablehnung eine eigene, unabhängige
Regimeeinschätzung abgeben. Ein fail-safe ausgelegter Validator prüft jede Antwort:
Schemaverletzungen, Zeitüberschreitungen, Widersprüche zwischen eigener Einschätzung
und Strategiesignal sowie Größenüberschreitungen führen zum Veto des betreffenden
Votums; ein nicht erreichbares Modell zählt konservativ als Ablehnung. Die
Einzelvoten werden nach strikter Mehrheit aggregiert, wobei bei Zustimmung die
konservativste (kleinste) Positionsgröße gilt. Steht eine ablehnende Mehrheit
vorzeitig fest, bricht die Befragung ab (\emph{Early-Stop}); eine bereits erreichte
Zustimmungsmehrheit führt hingegen bewusst nicht zum Abbruch, da weitere Voten die
Größe nur verkleinern können. Sämtliche Voten samt Begründungen und Antwortzeiten
werden im Entscheidungsjournal persistiert.

Die Orderausführung ist mehrstufig gegen die typischen Fehlerbilder verteilter
Systeme abgesichert. Ein Frische-Gate verwirft Entscheide, die nicht vom aktuellen
Handelstag stammen. Vor jeder realen Order wird eine \emph{Intent}-Zeile in das
Orderjournal geschrieben, die als eindeutiger Ausführungsmarker dient: Derselbe
Entscheid kann dadurch -- etwa nach einem manuellen Wiederholungslauf oder einer
doppelten Zeitplanung -- kein zweites Mal ausgeführt werden, und selbst ein
Prozessabsturz zwischen Orderaufgabe und Journalisierung bleibt nachvollziehbar.
Scheitert die Kommunikation mit der Bank nach dem Absenden, gleicht eine
Reconciliation den tatsächlichen Zustand über das Orderbuch der Bank ab; jeder
unklare Ausgang endet in einem alarmierten Handelsstopp statt in einem automatischen
Wiederholungsversuch. Einzelorder- und Portfolio-Obergrenzen wirken als
Circuit-Breaker, und die Scharfschaltung auf reale Orders ist ausschließlich
kontospezifisch möglich. Die Bankanbindung selbst nutzt den OAuth-Flow der
comdirect-Schnittstelle mit einer täglich per photoTAN freigegebenen Session-TAN;
der Mensch bleibt über einen Signal-Messenger-Kanal eingebunden, der Tagesreports
liefert und einen jederzeit wirksamen Not-Aus (Pausieren, Fortsetzen, Statusabfrage,
Glattstellen) bereitstellt.

Die Banksitzung selbst folgt einem täglichen Ritual, das die regulatorische
Zwei-Faktor-Anforderung in den autonomen Ablauf integriert: Der morgendliche
Anmeldevorgang fordert genau eine photoTAN-Freigabe an, die der Betreiber
fernbedient über den Messaging-Kanal bestätigt; anschließend hält ein
Erneuerungsdienst die Sitzung über den Handelstag, sodass Orders ohne weitere
Einzelfreigaben, aber stets innerhalb der freigegebenen Sitzung erfolgen. Die
gesamte Broker-Schicht ist zudem kontofähig ausgelegt: Zugangsdaten, Sitzungen,
Journale und der Ausführungsmodus sind je Konto getrennt, sodass der geplante
Parallelbetrieb zweier Depots mit unterschiedlichen Strategien keine strukturelle
Änderung erfordert. Zugangsdaten liegen verschlüsselt in einem dedizierten
Geheimnisspeicher und erscheinen zu keinem Zeitpunkt im Kontext der Sprachmodelle.

Das öffentliche Dashboard schließlich ist als eigenständiges, bewusst kleines
Teilsystem implementiert. Ein Exportprozess redigiert das Entscheidungsjournal auf
die veröffentlichungsfähigen Inhalte, lädt sie nur bei inhaltlicher Änderung hoch
und trennt so den öffentlichen Lesezugriff vollständig vom Handelspfad. Die
Oberfläche stellt den Indexverlauf mit den gleitenden Durchschnitten und den
Regimewechsel-Signalen dar, zeigt je Entscheidung das aggregierte Votum mit
aufklappbaren Einzelbegründungen der Modelle und ist durchgängig responsiv
umgesetzt -- auf Mobilgeräten werden Ordertabellen in vertikale
Beschriftungs-Wert-Karten umgebrochen, sodass die Darstellung ohne horizontales
Scrollen barrierearm lesbar bleibt.

Die Betriebsführung rundet die Implementierung ab. Ein Zeitplandienst startet den
Entscheidungszyklus an Handelstagen kurz vor der US-Markteröffnung; ein
Börsenkalender überspringt Feiertage, und eine Prozesssperre je Konto verhindert,
dass sich zwei Läufe überlappen. Der Not-Aus des Messaging-Kanals ist als
Ein-Buchstaben-Protokoll ausgeführt: Pausieren stoppt jede neue Order vor dem
nächsten Zyklus, Fortsetzen hebt die Sperre wieder auf, die Statusabfrage meldet
den Betriebszustand samt wirksamem Ausführungsmodus, und das Glattstellen schließt
offene Positionen aktiv -- wobei die Pause bewusst erhalten bleibt, bis der Mensch
sie aufhebt. Jeder Zyklus schließt mit einem Bericht an die private
Messaging-Gruppe, der Entscheidung, Voten, Auffälligkeiten und den effektiven
Modus je Konto ausweist; die Scharfschaltung auf reale Orders ist damit täglich
sichtbar dokumentiert.

Die Qualitätssicherung stützt sich auf zwölf Offline-Testsuiten, die von der
Datensammlung über den Voting-Mechanismus bis zu den Schutzpfaden der
Orderausführung jede sicherheitsrelevante Verzweigung mit gemockten Schnittstellen
durchspielen; die Volltexte der zentralen Module sind im Anhang dokumentiert.

\section{Evaluation}
\label{sec:umsetzung-evaluation}

Die Bewertung des Agenten erfolgt auf drei Ebenen: einem historischen Backtest der
Strategie, einer Analyse der realen Transaktionskosten sowie einer laufenden
Kalibrierung des Sprachmodell-Votings im Betrieb.

\subsection{Untersuchungsdesign und Datenbasis}
\label{sec:umsetzung-evaluation-design}

Datenbasis des Backtests sind zehn Jahre echter Tagesschlusskurse des S\&P~500, die
über die Datensammlung des Agenten bezogen und gegen die FRED-Zweitquelle
gegengeprüft wurden. Der Strategie werden drei Vergleichsmaßstäbe gegenübergestellt:
der tägliche Round-Trip der Regel selbst (B0), das passive Halten des Index (B1)
sowie eine Variante, die nur bei einem Regimewechsel handelt (B2). Das Kostenmodell
verwendet keine Annahmen, sondern die real gemessenen Ex-Ante-Kostenausweise der
Bank -- im Direkthandel 7,80\,€ je Round-Trip -- und berücksichtigt die
Kapitalertragsteuer unter Teilfreistellung. Ergänzend misst ein wöchentlicher
Auswertungslauf das Abstimmungsverhalten der Sprachmodelle im laufenden Betrieb: je
Modell die Zustimmungsquote, die Schemafehlerrate, die Antwortlatenz sowie die
Veto-Trefferquote gegen den Folgetags-Return, also den Anteil der Ablehnungen, die
tatsächlich einen Verlusttag in Signalrichtung vermieden hätten, einschließlich des
kontrafaktischen Gewinn- und Verlustbeitrags der Vetos.

Zur Einordnung der Handelsfrequenz dient zusätzlich ein einfaches
Rentabilitätsmodell. Bezeichnet $C(S)$ die Round-Trip-Kosten bei Kapitaleinsatz $S$
und wird die mittlere absolute Intraday-Bewegung des Index mit rund $0{,}6\,\%$
angesetzt, so ergibt sich die für einen positiven Erwartungswert erforderliche
Trefferquote als
\begin{equation}
  p^{*}(S) \;=\; 0{,}5 \;+\; \frac{C(S)}{4 \cdot S \cdot 0{,}006}.
  \label{eq:breakeven}
\end{equation}
Für 1.000\,€ Einsatz liegt diese Schwelle bei etwa 91\,\%, für 2.000\,€ bei etwa
75\,\%; selbst im Grenzfall sehr großer Einsätze setzt der Spread eine
kapitalunabhängige Untergrenze von rund 58\,\% -- Werte, denen die Literatur für
Intraday-Signale einfacher Durchschnittsstrategien lediglich Trefferquoten nahe der
Zufallsschwelle gegenüberstellt. Das Modell begründet damit vorab, weshalb die
tägliche Handelsfrequenz als Demonstrations- und nicht als Renditevehikel zu
verstehen ist.

\subsection{Ergebnisbetrachtung}
\label{sec:umsetzung-evaluation-ergebnisse}

Das zentrale Backtest-Ergebnis fällt eindeutig aus: Der tägliche Round-Trip (B0)
verliert über zehn Jahre 29.598,79\,€, wovon 28.937,94\,€ auf Transaktionskosten
entfallen, bei einer Trefferquote von 31,7\,\%. Die nur bei Regimewechseln handelnde
Variante (B2) erzielt mit neun Positionswechseln in zehn Jahren rund $+100\,\%$ bei
lediglich 106\,€ Kosten; das passive Halten des Index (B1) rund $+250\,\%$. Die
tägliche Handelsfrequenz wird demnach von den Fixkosten aufgezehrt -- die
Rentabilitätsanalyse zeigt, dass bei realistischen Trefferquoten kein
Kapitaleinsatz die tägliche Round-Trip-Strategie in den positiven Erwartungswert
hebt. Diese Erwartung wurde bewusst ehrlich formuliert: Das Erkenntnisziel der
Arbeit ist nicht eine positive Handelsrendite, sondern die autonome, abgesicherte
Ausführung selbst sowie der messbare Beitrag des Sprachmodell-Gremiums; die
Differenz zwischen B0 und B2 beziffert insofern den Preis der täglich demonstrierten
Autonomie. Die Voting-Kalibrierung aus dem Live-Betrieb -- Trockenlauf seit dem
21.07.2026, Echtgeldbetrieb ab dem 27.07.2026 -- ergänzt diese Betrachtung
fortlaufend um die Frage, ob die Vetos der Modelle systematisch Verlusttage treffen
oder lediglich Rauschen darstellen.

Die Kostenmessung gegen das Produktivsystem der Bank lieferte darüber hinaus
eigenständige Befunde. Die Kaufprovision erwies sich über alle getesteten
Ordergrößen als fixer Betrag von 3{,}90\,€, womit sich die Round-Trip-Fixkosten von
7{,}80\,€ bestätigen; die Wahl des Direkthandels beim Emittenten statt eines
Börsenplatzes senkt die Gesamtkosten je Round-Trip um weitere 5\,€, da die
Handelsplatzentgelte entfallen. Zugleich zeigte sich ein versteckter Kostenfaktor:
Der Geld-Brief-Spread des Direkthandels wird im Ex-Ante-Kostenausweis nicht
ausgewiesen und muss im Betrieb empirisch gemessen werden. Auch der Betrieb der
Sprachmodelle erbrachte belastbare Beobachtungen: Die Schema-Adhärenz kleinerer
gehosteter Modelle musste schrittweise erzwungen werden -- von anfänglich
verworfenen Antworten bis zur stabilen, schema-konformen Ausgabe --, für
Modelle mit aktivierter Reasoning-Ausgabe war deren Abschaltung erforderlich, weil
der Denkmodus das erzwungene Antwortformat überstimmte, und ein zunächst
vorgesehenes viertes Modell wurde nach Verifikationsläufen wegen langer
Antwortzeiten und inhaltlichen Rauschens bewusst ausgeschlossen. Die
Antwortlatenzen fielen nach der Warmlaufphase von knapp zehn Sekunden auf unter
eine Sekunde je Votum.

\subsection{Limitationen}
\label{sec:umsetzung-evaluation-limitationen}

Die Aussagekraft der Ergebnisse unterliegt mehreren Einschränkungen. Die
Bankschnittstelle bietet keine Sandbox, sodass die Ordervalidierung gegen das
Produktivsystem erfolgt und die Ausführungsseite bis zum Echtgeldbetrieb nur bis zur
Kostenvorschau demonstriert werden konnte. Die Autonomie ist regulatorisch durch die
täglich erforderliche photoTAN-Freigabe der Banksitzung begrenzt. Im verteilten
Aufbau beträgt die Datenfrische am Entscheidungspunkt 60 bis 90 Sekunden, weshalb
ausschließlich Limit-Orders verwendet werden. Der Backtest approximiert das
Faktor-Zertifikat über das zweifache Index-Tagessegment und modelliert weder
Währungseffekte noch den variablen Geld-Brief-Spread exakt; die inoffizielle
Kursdatenquelle bietet keine Verfügbarkeitsgarantie. Trotz Temperatur null und
erzwungenem Antwortschema bleibt ein Restmaß an Nichtdeterminismus der
Sprachmodelle bestehen, und die Stichprobe der Voting-Kalibrierung ist zu Beginn
des Betriebs klein, weshalb die Auswertung bewusst deskriptiv bleibt und keine
automatische Modellabschaltung auslöst.

\section{Kritische Würdigung}
\label{sec:umsetzung-wuerdigung}

Autonomer Handel mit gehebelten Derivaten vereint mehrere Risikoklassen: das
Marktrisiko des gehebelten Instruments, das Betriebsrisiko eines unbeaufsichtigt
laufenden Systems und das spezifische Risiko probabilistischer Modellkomponenten.
Die Architektur begegnet dem mit einer gestaffelten Sicherungskette, deren
tragendes Prinzip die begrenzte Diskretion ist: Das Sprachmodell-Gremium kann
Handelsvorschläge verhindern, verkleinern oder zeitlich einordnen, niemals aber
erzeugen oder in ihrer Richtung umkehren, und deterministische Ausschlusskriterien
werden grundsätzlich vor jedem Modellaufruf entschieden. Auf der Ausführungsseite
verhindern das zweiphasige Intent-Journal die Doppelausführung, die
Orderbuch-Reconciliation unbemerkte Zustandslücken nach Kommunikationsfehlern und
die Einzel- sowie Portfolio-Obergrenzen ein unkontrolliertes Anwachsen der
Positionen; jeder unklare Zustand mündet in einem alarmierten Handelsstopp. Der
Mensch bleibt dreifach eingebunden: durch die tägliche photoTAN-Freigabe als
regulatorisch erzwungenen Kontrollpunkt, durch den in Echtzeit wirksamen Not-Aus
über den Messaging-Kanal und durch einen Totmannschalter, der den Ausfall des
Agenten selbst meldet. Das öffentliche Dashboard ergänzt die Kontrolle um
Transparenz, ohne die handelnden Systeme zu exponieren.

Gleichwohl bleiben Restrisiken bestehen. Der Ausstiegspfad bei einem Regimewechsel
ist bis zum geplanten Soll-Positions-Abgleich bewusst auf einen alarmierten Stopp
mit manueller Behandlung reduziert -- eine konservative Zwischenlösung, die
menschliches Eingreifen erfordert, statt fehlerhaft zu automatisieren. Die
Wirksamkeit des Modell-Gremiums ist zu Betriebsbeginn empirisch noch unbelegt und
wird erst durch die laufende Kalibrierung überprüfbar; korrelierte Fehlurteile der
Modelle bleiben trotz Anbieterdiversität möglich, da alle Modelle denselben
Entscheidungskontext erhalten. Wirtschaftlich ist festzuhalten, dass die täglich
handelnde Demonstrationsstrategie nach Kosten strukturell defizitär ist und der
Agent insofern ein Erkenntnis-, kein Renditeinstrument darstellt. Von einer
Anlageberatung grenzt sich die Arbeit ausdrücklich ab: Der Agent handelt
ausschließlich für das eigene Experimentalkonto des Betreibers, und weder die
Strategie noch die Systemarchitektur stellen eine Empfehlung an Dritte dar.
Schließlich zeigte die Umsetzung einen strukturellen Befund des Marktumfelds: Die
kostengünstigsten Ausführungswege für Privatanleger sind systematisch nicht über
offene Schnittstellen zugänglich, sodass autonome Agenten einen faktischen
\enquote{API-Aufpreis} tragen -- ein Umstand, der die angestrebte, weitgehend
kostenfreie lokale Betreibbarkeit umso relevanter macht, da sie zumindest die
Infrastruktur- und Inferenzkosten eliminiert.

Über den konkreten Anwendungsfall hinaus beansprucht die Architektur
Übertragbarkeit. Strategie, Produktzuordnung, Modellgremium und Bankadapter sind
über definierte Kontrakte entkoppelt, sodass sich einzelne Bausteine austauschen
lassen, ohne die Sicherungskette zu berühren: Eine andere Regelstrategie ist ein
weiteres Plugin, ein anderes Sprachmodell eine Konfigurationszeile, ein anderer
Broker ein neuer Adapter hinter derselben Orderabstraktion. Für die geplante
Migration auf einen Mac~mini verbleiben als offene Arbeiten im Kern die Anpassung
der Zeitbudgets an langsamere lokale Inferenz, der Ersatz einzelner
Linux-spezifischer Betriebswerkzeuge und die Umstellung der Zeitplanung -- der
Python-Kern selbst ist frei von Cloud-Abhängigkeiten gehalten. Damit bleibt das
Versprechen der Zielsetzung, dieselbe Architektur vom Cloud-Prototyp in einen
privat betreibbaren, laufend kostenfreien Betrieb zu überführen, strukturell
eingelöst, auch wenn der empirische Nachweis dieser letzten Stufe aussteht.

%  =============================================================================
%  Ende
%  =============================================================================

%  =============================================================================
%  Datei-Hinweise & Konfiguration: siehe config.md
%  -> Abschnitt "kapitel/04_schluss.tex"
%  =============================================================================

\chapter{Schluss}
\label{cha:schluss}

Nachdem der Handelsagent entworfen, implementiert und evaluiert wurde, fasst dieses
Kapitel die zentralen Erkenntnisse mit Blick auf die Zielsetzung zusammen, diskutiert
Stärken und Schwächen der gewählten Lösung und gibt einen Ausblick auf
weiterführende Arbeiten.

\section{Fazit}
\label{sec:fazit}

Ziel dieser Arbeit war der Aufbau eines Live-Trading-Agenten auf Basis des
Agenten-Frameworks OpenClaw, dessen deterministische \glsxtrshort{sma}-Strategie
durch ein Gremium großer Sprachmodelle geprüft wird, bevor ein Mehrheitsentscheid
die Orderausführung freigibt (\cref{sec:zielsetzung}). Dieses Ziel wurde erreicht:
Der vollständige Handelszyklus -- von der kontinuierlichen, gegen eine
Referenzreihe abgesicherten Marktdatenerfassung über die regelbasierte
Regimeentscheidung und das Multi-Modell-Voting bis zur Orderaufgabe über die
Schnittstelle einer regulierten Bank -- ist umgesetzt, im Trockenlauf erprobt und
in \cref{cha:umsetzung} vollständig dokumentiert. Die theoretischen Grundlagen aus
\cref{cha:grundlagen} bildeten dafür den methodischen Rahmen: Das Periodensystem
der Künstlichen Intelligenz erlaubte es, die bewusste Begrenzung des
\glsxtrshort{ki}-Einsatzes nicht als Verzicht, sondern als gestalterische Aussage
zu formulieren -- die Elemente der Kursprognose und der selbstlernenden
Fortschreibung bleiben unbesetzt, das Element der prozesssteuernden Kontrolle
bildet das Zentrum.

Mit Blick auf die erste Leitfrage zeigt die Fallstudie, dass der autonome,
nachvollziehbare Betrieb mit vertretbarem Aufwand möglich ist -- der wesentliche
Aufwand liegt dabei nicht im Handelsweg, sondern in seiner Absicherung: im
zweiphasigen Ausführungsjournal gegen Doppelorders, im Orderbuchabgleich nach
Kommunikationsfehlern, in den Kapitalobergrenzen sowie im jederzeit wirksamen
Not-Aus (\cref{sec:umsetzung-implementierung}). Die zweite Leitfrage nach dem
messbaren Beitrag des Mehrheitsvotums bleibt zum Abgabezeitpunkt bewusst offen;
die dafür geschaffene Kalibrierung misst den Veto-Beitrag der Modelle fortlaufend
gegen die tatsächliche Marktentwicklung und benötigt eine längere Betriebsdauer,
bevor valide Aussagen möglich sind (\cref{sec:umsetzung-evaluation-design}).

In Kennzahlen verdichtet stellt sich das Ergebnis wie folgt dar: Der Backtest
über zehn Jahre stellt dem täglichen Round-Trip der Demonstrationsstrategie
(rund $-29.700$\,€, davon etwa 98\,\% Transaktionskosten) das reine Regimefolgen
(rund $+96\,\%$ vor Steuern bei neun Positionswechseln) und das passive Halten
(rund $+248\,\%$ vor Steuern) gegenüber; die Walk-Forward-Prüfung des
Prognosemodells über
23.803 Handelstage ergibt keine Prognosekraft über die Basisrate hinaus bei
zugleich guter Kalibrierung; die real gemessenen Fixkosten betragen 7,80\,€ je
Round-Trip, und die Antwortlatenzen des Modellgremiums liegen im eingeschwungenen
Betrieb unter einer Sekunde je Votum
(\cref{sec:umsetzung-evaluation-ergebnisse}). Ökonomisch fällt die Antwort damit
eindeutig aus: Die täglich handelnde Demonstrationsstrategie ist nach realen
Kosten strukturell defizitär; der Agent ist ein Erkenntnis-, kein
Renditeinstrument, und die Differenz zum reinen Regimefolgen beziffert den Preis
der täglich demonstrierten Autonomie. Methodisch hat sich dabei das iterative
Vorgehen bewährt, jede Ausbaustufe unmittelbar gegen reale Betriebsdaten zu
prüfen: Mehrere der wirksamsten Absicherungen -- von der Unterscheidung legitimer
Datenstille von echten Verbindungsverlusten bis zur datumskorrekten
Zweitquellenprüfung -- entstanden nicht am Reißbrett, sondern aus Befunden des
laufenden Betriebs und wurden anschließend durch Regressionstests dauerhaft
abgesichert.

\section{Diskussion}
\label{sec:diskussion}

Den Stärken der gewählten Lösung stehen benennbare Schwächen gegenüber; beide
werden hier bewusst gebündelt, um die Zielerreichung ehrlich einzuordnen. Als
tragfähig erwiesen hat sich erstens das Prinzip der begrenzten Diskretion: Kein
beobachteter Fehlerfall -- weder Schemaverletzungen einzelner Modelle noch
Zeitüberschreitungen oder Datenlücken -- konnte eine unbeabsichtigte Order
auslösen, weil jede Unklarheit konservativ als Veto oder Handelsstopp endet.
Zweitens hat sich die Entscheidung ausgezahlt, die Evaluation auf gemessene statt
angenommene Größen zu stützen: Reale Kostenausweise, die laufende
Kursgegenprüfung und die Walk-Forward-Prüfung machen die Ergebnisse belastbar, und der ehrliche
Negativbefund zur Prognosekraft trägt die zentrale Architekturentscheidung
empirisch, statt sie zu beschädigen. Drittens schafft die durchgängige
Journalisierung zusammen mit dem öffentlichen Dashboard eine Transparenz, die
jede einzelne Entscheidung samt Modellvoten nachvollziehbar und den laufenden
Betrieb öffentlich prüfbar macht.

Dem stehen Schwächen gegenüber. Der erhoffte Sicherheitsgewinn des
Modellgremiums ist zu Betriebsbeginn empirisch unbelegt, und korrelierte
Fehlurteile bleiben trotz Anbieterdiversität möglich, da alle Modelle denselben
Entscheidungskontext erhalten; die Kalibrierung kann dies erst nach längerer
Betriebsdauer klären. Die tägliche Handelsfrequenz ist nach Kosten strukturell
defizitär und nur als Demonstrationsvehikel zu rechtfertigen. Die Autonomie
bleibt an regulatorische und betriebliche Grenzen gebunden -- die tägliche
photoTAN-Freigabe, den bewusst manuell gehaltenen Ausstiegspfad bei
Regimewechseln und eine Datenfrische von 60 bis 90 Sekunden am
Entscheidungspunkt. Schließlich stützt sich die Kursdatenerfassung auf eine
inoffizielle Quelle ohne Verfügbarkeitsgarantie, deren Ausfall der Betrieb zwar
erkennen, aber nicht kompensieren kann. Aus dieser Gegenüberstellung folgt das
Verbesserungspotenzial, das der nachstehende Ausblick in konkrete
Weiterentwicklungen übersetzt.

\section{Ausblick}
\label{sec:ausblick}

Aus der Bearbeitung haben sich die folgenden weiterführenden Ansatzpunkte ergeben,
die Gegenstand künftiger Arbeiten sein können. Als Erstes zu nennen ist der
unmittelbar bevorstehende Übergang vom Trockenlauf in den Echtgeldbetrieb, der die
bislang bis zur Ordervalidierung demonstrierte Ausführungsseite unter realen
Bedingungen prüft und die Datenbasis der Voting-Kalibrierung kontinuierlich
erweitert. Als weiterer Punkt kann der Soll-Positions-Abgleich umgesetzt werden,
der den Ausstiegspfad bei Regimewechseln vollständig deterministisch macht und
damit die letzte bewusst manuell gehaltene Zwischenlösung ablöst; darauf aufbauend
erlaubt der Parallelbetrieb eines zweiten Depots den direkten Vergleich der täglich
handelnden Kreuzungsstrategie mit der haltenden Regimestrategie unter identischen
Bedingungen. Des Weiteren bietet sich die Weiterentwicklung des Modellgremiums hin
zu komplementären Prüfrollen an, um der Gefahr korrelierter Fehlurteile zu
begegnen; die laufende Kalibrierung liefert hierfür die datenbasierte
Entscheidungsgrundlage, welche Modelle das Gremium tatsächlich stärken.
Abschließend wäre die Migration auf einen Mac~mini mit lokal betriebenen Modellen
anzustreben, die die laufenden Kosten weitgehend eliminiert und die Datenhoheit
vervollständigt -- die Architektur ist darauf vorbereitet, der Wechsel bleibt eine
Konfigurationsentscheidung.

Damit liegt mit dieser Arbeit eine nachvollziehbar dokumentierte, übertragbare
Referenzarchitektur für sicherheitsorientiert begrenzten \glsxtrshort{ki}-Einsatz
im autonomen Handel vor, die als Grundlage für die weiterführende Untersuchung der
Frage dient, unter welchen Bedingungen probabilistische Sprachmodelle als
kontrollierte Prüfinstanzen deterministischer Systeme einen messbaren
Sicherheitsgewinn stiften.

%  =============================================================================
%  Ende
%  =============================================================================


% --- Nachspann --------------------------------------------------------------
% Leitfaden Kap. 5.1 / 6.3: Quellenverzeichnis nach dem Hauptteil,
% alphabetisch nach Erstautor. Kap. 6.2.4: KI-Verzeichnis folgt direkt
% nach dem Literaturverzeichnis.
\cleardoublepage
\printbibliography[
    filter=nichtKI,
    title={Literaturverzeichnis},
    heading=bibintoc
]

% KI-Verzeichnis (Kap. 6.2.4) -- ausgelagert in eigene Verzeichnis-Datei,
% direkt hinter dem Literaturverzeichnis eingebunden.
\cleardoublepage
%  =============================================================================
%  Datei-Hinweise & Konfiguration: siehe config.md
%  -> Abschnitt "verzeichnisse/ki-verzeichnis.tex"
%  =============================================================================

\printbibliography[
    filter=nurKI,
    title={KI-Verzeichnis},                      % Kap. 6.2.4
    heading=bibintoc
]

%  =============================================================================
%  Ende
%  =============================================================================


% Anhang mit grossen roemischen Ziffern auf ALLEN Gliederungsebenen
% (Kap. 5.1: "Die Nummerierung im Anhang erfolgt mit grossen roemischen
% Ziffern.")
\appendix
\renewcommand{\thechapter}{\Roman{chapter}}
\renewcommand{\thesection}{\thechapter.\Roman{section}}
\renewcommand{\thesubsection}{\thesection.\Roman{subsection}}
\renewcommand{\thesubsubsection}{\thesubsection.\Roman{subsubsection}}
% =============================================================================
%  Datei-Hinweise & Konfiguration: siehe config.md
%  -> Abschnitt "anhang/anhang.tex"
% =============================================================================

% Anhang I: Projektstruktur + Quellcodearchiv (Verweisziel aus Kapitel 3;
%           ersetzt seit V1_8 den Volltext-Abdruck anhang/quelltexte.tex),
% Anhang II: Evaluationsbelege (Verweisziel aus Kapitel 3.4),
% Anhang III: KI-Prompts (Kap. 6.2.4).
%  =============================================================================
%  Datei-Hinweise & Konfiguration: siehe config.md
%  -> Abschnitt "anhang/quellcodearchiv.tex"
%
%  PFLEGE / REGENERIERUNG (z. B. nach der Echtgeld-Woche):
%    1. Im Projektordner archiv_manifest.csv anpassen (neue Datei = neue
%       Zeile; Status "geplant" -> "aktiv" setzen).
%    2. python3 build_quellcode_archiv.py --version V3 --datum <TTMONJJJJ>
%    3. Die vier GENERIERTEN Dateien archiv_meta.tex, archiv_quellcode.tex,
%       archiv_tests.tex, archiv_daten.tex aus archiv_tex/ nach Overleaf in
%       den Ordner anhang/ hochladen (ersetzen). Diese Datei hier bleibt
%       unveraendert.
%  =============================================================================

\chapter{Anhang: Projektstruktur und Quellcodearchiv}
\label{cha:anhang-archiv}

% Generierte Metadaten des Archivs (Name, Datum, SHA-256, Umfang).
%% ---------------------------------------------------------------------------
%% GENERIERT durch build_quellcode_archiv.py am 27.07.2026 12:00 — NICHT von Hand
%% pflegen; Aenderungen in archiv_manifest.csv vornehmen und Skript ausfuehren.
%% ---------------------------------------------------------------------------
\newcommand{\QuellcodeArchivName}{FOM\_Projektarbeit\_KI\_Quellcode\_FischerKiss\_27JUL2026\_V3.zip}
\newcommand{\QuellcodeArchivDatum}{27.07.2026}
\newcommand{\QuellcodeArchivVersion}{V3}
\newcommand{\QuellcodeArchivSHA}{2c56790b80c2f793fa47665b041840a27aefa26596cb04ea80d4560772d5a09e}
\newcommand{\QuellcodeArchivDateien}{112}
\newcommand{\QuellcodeArchivGroesse}{0.7}


Dieser Anhang dokumentiert den Quellcodestand des in \cref{cha:umsetzung}
beschriebenen Handelsagenten. Der vollständige Quellcode wird nicht im Anhang
abgedruckt, sondern liegt der Arbeit als elektronisches Archiv bei; die
nachfolgenden Abschnitte weisen dessen Identität, Aufbau und Inhalt aus. Damit
bleibt jede in \cref{sec:umsetzung-implementierung} beschriebene Absicherung im
Quelltext nachvollziehbar, ohne den Anhang mit rund 7.000 Zeilen Quelltext zu
belasten.

\section{Projektstruktur}
\label{sec:anhang-projektstruktur}

\Cref{fig:projektstruktur} zeigt die Struktur des Projekts. Es handelt sich um
einen einzigen Codestand, der auf beiden virtuellen Maschinen ausgerollt wird;
die Zuordnung der Module zu collector-vm und agent-vm ist je Datei vermerkt.
Die Zuordnung jeder Datei zu ihrem Zweck und ihrem Ort im Archiv dokumentieren
die Tabellen der \cref{sec:anhang-archiv-quellcode,sec:anhang-archiv-tests,sec:anhang-archiv-daten}.

\begin{figure}[htbp]
    \centering
    \includesvg[width=0.92\textwidth,height=0.88\textheight,keepaspectratio]{projektstruktur}
    \caption[Projektstruktur des Handelsagenten]{Projektstruktur des
        Handelsagenten: ein Codestand für beide virtuellen Maschinen
        (C = collector-vm, A = agent-vm).\quelleeigenki{Anthropic Claude Fable~5 (Anthropic, 2026)}}
    \label{fig:projektstruktur}
\end{figure}

\section{Archividentität und Redaktion}
\label{sec:anhang-archiv-identitaet}

Das Archiv ist über Dateinamen, Erstellungsdatum und Prüfsumme eindeutig
identifiziert; die Prüfsumme erlaubt den Nachweis, dass der beigelegte Stand
dem hier dokumentierten entspricht.

\begin{table}[htbp]
  \centering
  \begin{tabular}{l l}
    \hline
    Merkmal & Wert \\
    \hline
    Dateiname & \texttt{\footnotesize\QuellcodeArchivName} \\
    Version / Datum & \QuellcodeArchivVersion{} / \QuellcodeArchivDatum \\
    Umfang & \QuellcodeArchivDateien{} Dateien, \QuellcodeArchivGroesse{}\,MB \\
    Integrität (SHA-256) & \texttt{\footnotesize\QuellcodeArchivSHA} \\
    \hline
  \end{tabular}
  \caption[Identität des Quellcodearchivs]{Identität des beigelegten
    Quellcodearchivs.\quelleeigen}
  \label{tab:anhang-archiv-identitaet}
\end{table}

Für die Beilage gelten zwei redaktionelle Regeln, die die Programmlogik
unberührt lassen: Die Rufnummern des Messaging-Kanals sind redigiert
(\texttt{+49XXXXXXXXXXX}), und Zugangsdaten erscheinen an keiner Stelle, da sie
ausschließlich im verschlüsselten Geheimnisspeicher liegen -- die Datei
\texttt{agent.env} und der Inhalt von \texttt{\textasciitilde/.fom-agent/}
sind deshalb nicht Bestandteil des Archivs. Die unverändert eingebundene
Fremdbibliothek Chart.js (\texttt{webfeed/app/static/chart.umd.js}, MIT-Lizenz)
ist als Drittcode ebenfalls nicht beigelegt; im Übrigen liegt der Quellcode
unverändert und lauffähig vor.

\section{Aufbau des Archivs}
\label{sec:anhang-archiv-aufbau}

Das Archiv gliedert sich in drei Basisordner sowie eine erläuternde
\texttt{README.md}: Der Ordner \texttt{Quellcode/} enthält sämtliche Module,
Skripte und Konfigurationsdateien in der flachen Struktur der
Produktivumgebung (vgl. \cref{fig:projektstruktur}), wodurch der Code ohne
Umbau lauffähig bleibt; der Ordner \texttt{Tests/} enthält die
Offline-Testsuiten einschließlich des Testtreibers \texttt{run\_tests.sh}; der
Ordner \texttt{Daten/} enthält die Daten- und Evaluationsstände mit
ausgewiesenem Snapshot-Datum.

\section{Quellcode}
\label{sec:anhang-archiv-quellcode}

%% ---------------------------------------------------------------------------
%% GENERIERT durch build_quellcode_archiv.py am 25.07.2026 08:59 — NICHT von Hand
%% pflegen; Aenderungen in archiv_manifest.csv vornehmen und Skript ausfuehren.
%% ---------------------------------------------------------------------------
\begin{longtable}{@{}p{0.29\textwidth}p{0.155\textwidth}p{0.455\textwidth}@{}}
  \caption[Quellcodedateien des Archivs]{Ordner \texttt{Quellcode/} des Archivs: Dateien und Zwecke, gegliedert nach Funktionsbereichen der Projektstruktur.\quelleeigen}\label{tab:anhang-archiv-quellcode}\\
  \toprule
  Datei & Typ & Zweck \\
  \midrule
  \endfirsthead
  \multicolumn{3}{@{}l}{\itshape Fortsetzung von vorheriger Seite}\\
  \toprule
  Datei & Typ & Zweck \\
  \midrule
  \endhead
  \midrule
  \multicolumn{3}{r@{}}{\itshape Fortsetzung auf der n\"achsten Seite}\\
  \endfoot
  \bottomrule
  \endlastfoot
  \multicolumn{3}{@{}l}{\textbf{Konfiguration}}\\*[2pt]
  \texttt{agentconfig.yaml} & Kon\-fi\-gu\-ra\-tion & Betriebskonfiguration: Strategien, Deployments, Sprachmodelle, Zeitbudgets, ML-Prüfinstanz \\
  \texttt{wknassign.yaml} & Kon\-fi\-gu\-ra\-tion & Zentrale Produktzuordnung Strategie zu Handelsinstrument (WKN/ISIN) \\
  \addlinespace[6pt]
  \multicolumn{3}{@{}l}{\textbf{Datenerfassung (collector-vm)}}\\*[2pt]
  \texttt{datacollect.py} & Python & Kontinuierliche Marktdatenerfassung mit Sanitätsfilter, Bar-Verdichtung und FRED-Zweitquellenprüfung \\
  \texttt{market\_\allowbreak calendar.py} & Python & Börsenkalender: Handelstage, Feiertage, Sitzungsfenster \\
  \texttt{collector\_\allowbreak summary.py} & Python & Tageszusammenfassung der Datenerfassung für den Berichtsversand \\
  \addlinespace[6pt]
  \multicolumn{3}{@{}l}{\textbf{Strategieschicht}}\\*[2pt]
  \texttt{strategy.py} & Python & Strategie-Handler: Signalberechnung, DecisionRequest, Audit-Flags, Kostenkontext, Instrumentenkatalog, Backtest-Kommando \\
  \texttt{strategyloader.py} & Python & Plugin-Lader: aktive Strategien, Deployments, ML-Konfiguration \\
  \texttt{plstrategy\_\allowbreak sma.py} & Python & Plugin der SMA-Regimestrategie (Kreuzungssignal) \\
  \texttt{plstrategy\_\allowbreak smacross.py} & Python & Plugin-Variante des Kreuzungssignals \\
  \texttt{plstrategy\_\allowbreak smatrend.py} & Python & Trendfolge-Alternative für den A/B-Vergleich \\
  \texttt{plstrategy\_\allowbreak buyhold.py} & Python & Vergleichsmaßstab passives Halten \\
  \texttt{plstrategy\_\allowbreak momentum.py} & Python & Vergleichsmaßstab Momentum \\
  \texttt{wknassign\_\allowbreak sma.py} & Python & Veralteter Platzhalter, ersetzt durch wknassign.yaml (dokumentiert die Migration) \\
  \texttt{wknassign\_\allowbreak smacross.py} & Python & Veralteter Platzhalter, ersetzt durch wknassign.yaml \\
  \texttt{wknassign\_\allowbreak momentum.py} & Python & Veralteter Platzhalter, ersetzt durch wknassign.yaml \\
  \addlinespace[6pt]
  \multicolumn{3}{@{}l}{\textbf{Prüfinstanzen}}\\*[2pt]
  \texttt{auditor.py} & Python & Multi-LLM-Voting: Pre-Gate, unabhängige Voten, strikte Mehrheit, fail-safe Validator \\
  \texttt{mlforecast.py} & Python & Statistische Prüfinstanz: logistische Regression, Walk-Forward-Prüfung, Schatten- und Wirkmodus \\
  \addlinespace[6pt]
  \multicolumn{3}{@{}l}{\textbf{Orderausführung und Bankanbindung}}\\*[2pt]
  \texttt{orchestrate.py} & Python & Abgesicherte Orderausführung: Frische-Gates, Intent-Journal, Orderbuch-Reconciliation, Circuit-Breaker \\
  \texttt{broker.py} & Python & comdirect-Adapter: OAuth2, Session-TAN, Depot- und Orderbuchzugriff, tagesgültige Limit-Orders \\
  \texttt{credstore.py} & Python & Verschlüsselte Ablage der Zugangsdaten, getrennt je Konto \\
  \texttt{config.py} & Python & Lader der Umgebungskonfiguration \\
  \addlinespace[6pt]
  \multicolumn{3}{@{}l}{\textbf{Betrieb und Steuerung}}\\*[2pt]
  \texttt{run\_\allowbreak cycle.sh} & Shell & Täglicher Entscheidungszyklus: Preflight, Strategie, Audit, Ausführung, Bericht \\
  \texttt{run\_\allowbreak all.sh} & Shell & Cron-Einstieg: Entscheidungszyklus je Konto/Deployment \\
  \texttt{control.py} & Python & Befehlsprotokoll des Steuerkanals: Not-Aus, Broker-Login, ML-Umschaltung, Status \\
  \texttt{signal\_\allowbreak dispatcher.py} & Python & Empfangsschleife des Messaging-Kanals mit Absender-Autorisierung \\
  \texttt{signal\_\allowbreak dispatcher.service} & Kon\-fi\-gu\-ra\-tion & systemd-Unit der Empfangsschleife \\
  \texttt{notify.py} & Python & Berichts- und Alarmversand mit Stale-Guard \\
  \texttt{notify\_\allowbreak tan.sh} & Shell & photoTAN-Benachrichtigung beim Broker-Login \\
  \texttt{alert\_\allowbreak relay.sh} & Shell & Alarm-Weiterleitung der Datenerfassung an den Messaging-Kanal \\
  \texttt{cleanup\_\allowbreak groups.py} & Python & Werkzeug: Aufräumen verwaister Messaging-Gruppen \\
  \texttt{keepalive.sh} & Shell & Sitzungserhalt der Broker-Session (V2, Overnight-fähig) mit TAN-Fallback \\
  \texttt{keepalive.service} & Kon\-fi\-gu\-ra\-tion & systemd-Unit des Sitzungserhalts \\
  \texttt{keepalive.timer} & Kon\-fi\-gu\-ra\-tion & systemd-Timer des Sitzungserhalts \\
  \texttt{healthcheck.sh} & Shell & Systemüberwachung agent-vm: Dienste, Speicherplatz, Zyklusstatus \\
  \texttt{cycle\_\allowbreak watch.sh} & Shell & Wächter: Alarm bei ausbleibendem Zyklusergebnis \\
  \texttt{heartbeat.sh} & Shell & Heartbeat der Datenerfassung (Lebenszeichen-Datei) \\
  \texttt{heartbeat\_\allowbreak watch.sh} & Shell & Dead-Man-Überwachung des Heartbeats mit Alarmversand \\
  \texttt{loadprobe.sh} & Shell & Lastmessung der agent-vm (Dimensionierungsnachweis) \\
  \texttt{maintain.sh} & Shell & Wartungslauf der Datenerfassung (Rotation, Aufräumen) \\
  \texttt{vacuum.sh} & Shell & SQLite-Vacuum der Marktdatenbank \\
  \texttt{extract\_\allowbreak live\_\allowbreak config.sh} & Shell & Extraktion der Live-Betriebskonfiguration (systemd, Cron, Litestream) von den VMs \\
  \addlinespace[6pt]
  \multicolumn{3}{@{}l}{\textbf{Auswertung}}\\*[2pt]
  \texttt{evaluate\_\allowbreak votes.py} & Python & Wöchentliche Voting-Kalibrierung und ML-Auswertung gegen den Folgetags-Return \\
  \addlinespace[6pt]
  \multicolumn{3}{@{}l}{\textbf{Werkzeuge und Diagnose}}\\*[2pt]
  \texttt{cost\_\allowbreak probe.py} & Python & Ex-Ante-Kostenmessung der Ordervalidierung gegen das Banksystem \\
  \texttt{venue\_\allowbreak probe.py} & Python & Sondierung der Handelsplätze und Orderarten (comdirect) \\
  \texttt{dbg\_\allowbreak positions.py} & Python & Diagnose der Depotbestände \\
  \addlinespace[6pt]
  \multicolumn{3}{@{}l}{\textbf{Dashboard und Journal-Export (webfeed)}}\\*[2pt]
  \texttt{webfeed/\allowbreak export\_\allowbreak journal.py} & Python & Redigierter Journal-Export des Handelsjournals für das Dashboard (collector-vm) \\
  \texttt{webfeed/\allowbreak app/\allowbreak main.py} & Python & Öffentliches Dashboard (FastAPI, Cloud Run) \\
  \texttt{webfeed/\allowbreak app/\allowbreak templates/\allowbreak index.html} & Web & Dashboard-Oberfläche (Diagramme, Journalansicht) \\
  \texttt{webfeed/\allowbreak app/\allowbreak templates/\allowbreak doku.html} & Web & Dokumentationsseite des Dashboards \\
  \texttt{webfeed/\allowbreak app/\allowbreak static/\allowbreak kf-logo.svg} & Web & Logo-Grafik des Dashboards \\
  \texttt{webfeed/\allowbreak app/\allowbreak requirements.txt} & Kon\-fi\-gu\-ra\-tion & Python-Abhängigkeiten des Dashboards \\
  \texttt{webfeed/\allowbreak app/\allowbreak Dockerfile} & Kon\-fi\-gu\-ra\-tion & Container-Bauvorschrift des Dashboards (Cloud Run) \\
  \texttt{webfeed/\allowbreak app/\allowbreak .dockerignore} & Kon\-fi\-gu\-ra\-tion & Build-Ausschlüsse des Container-Images \\
\end{longtable}


\section{Tests}
\label{sec:anhang-archiv-tests}

Die Testsuiten lagen in der Produktivumgebung neben den geprüften Modulen und
sind im Archiv in den Ordner \texttt{Tests/} ausgegliedert. Der Testtreiber
\texttt{run\_tests.sh} stellt den Modulpfad her und führt alle Suiten offline
aus (\texttt{bash Tests/run\_tests.sh}); Suiten mit fehlenden
Drittabhängigkeiten werden dabei als übersprungen ausgewiesen.

%% ---------------------------------------------------------------------------
%% GENERIERT durch build_quellcode_archiv.py am 25.07.2026 08:59 — NICHT von Hand
%% pflegen; Aenderungen in archiv_manifest.csv vornehmen und Skript ausfuehren.
%% ---------------------------------------------------------------------------
\begin{longtable}{@{}p{0.26\textwidth}p{0.26\textwidth}p{0.38\textwidth}@{}}
  \caption[Testsuiten des Archivs]{Ordner \texttt{Tests/} des Archivs: Offline-Testsuiten mit gepr\"uftem Artefakt und Pr\"ufzweck.\quelleeigen}\label{tab:anhang-archiv-tests}\\
  \toprule
  Testdatei & Gepr\"uftes Artefakt & Zweck \\
  \midrule
  \endfirsthead
  \multicolumn{3}{@{}l}{\itshape Fortsetzung von vorheriger Seite}\\
  \toprule
  Testdatei & Gepr\"uftes Artefakt & Zweck \\
  \midrule
  \endhead
  \midrule
  \multicolumn{3}{r@{}}{\itshape Fortsetzung auf der n\"achsten Seite}\\
  \endfoot
  \bottomrule
  \endlastfoot
  \texttt{run\_\allowbreak tests.sh} & alle Testsuiten & Testtreiber: führt alle Suiten mit korrektem Modulpfad (PYTHONPATH) gegen Quellcode/ aus \\
  \texttt{test\_\allowbreak auditor.py} & \texttt{auditor.py} & Multi-LLM-Voting: Mehrheitslogik, Pre-Gate, fail-safe Verhalten (gemockte LLM-Antworten) \\
  \texttt{test\_\allowbreak brokersession.py} & \texttt{broker.py} & Broker-Adapter: Session-, OAuth- und Orderbuchlogik gegen Attrappen \\
  \texttt{test\_\allowbreak control.py} & \texttt{control.py} & Befehlsprotokoll des Steuerkanals einschließlich Absender-Autorisierung \\
  \texttt{test\_\allowbreak credstore.py} & \texttt{credstore.py} & Ver- und Entschlüsselung der Zugangsdatenablage \\
  \texttt{test\_\allowbreak crosscheck.py} & \texttt{datacollect.py} & Zweitquellenprüfung und Sanitätsfilter der Marktdatenerfassung \\
  \texttt{test\_\allowbreak deployments.py} & \texttt{strategyloader.py}, \texttt{agentconfig.yaml} & Konsistenz der Deployment- und Strategiekonfiguration \\
  \texttt{test\_\allowbreak dispatcher.py} & \texttt{signal\_\allowbreak dispatcher.py} & Empfangsschleife des Messaging-Kanals \\
  \texttt{test\_\allowbreak evaluate\_\allowbreak votes.py} & \texttt{evaluate\_\allowbreak votes.py} & Voting- und ML-Auswertung gegen den Folgetags-Return \\
  \texttt{test\_\allowbreak keepalive.sh} & \texttt{keepalive.sh} & Sitzungserhalt: Szenarien einschließlich TAN-Fallback und Nachtruhe \\
  \texttt{test\_\allowbreak mlforecast.py} & \texttt{mlforecast.py} & Statistische Prüfinstanz: Merkmalsbildung, Walk-Forward, Schwellenlogik \\
  \texttt{test\_\allowbreak notify.py} & \texttt{notify.py} & Berichtsversand und Stale-Guard \\
  \texttt{test\_\allowbreak orchestrate.py} & \texttt{orchestrate.py} & Orderausführung: Frische-Gates, Intent-Journal, Reconciliation, Circuit-Breaker \\
  \texttt{webfeed/\allowbreak test\_\allowbreak export.py} & \texttt{webfeed/\allowbreak export\_\allowbreak journal.py} & Journal-Export: Redaktion und Aggregation \\
  \texttt{webfeed/\allowbreak app/\allowbreak test\_\allowbreak app.py} & \texttt{webfeed/\allowbreak app/\allowbreak main.py} & Dashboard-Endpunkte (FastAPI-Testclient) \\
\end{longtable}


\section{Daten}
\label{sec:anhang-archiv-daten}

Die Datenstände belegen die in \cref{sec:umsetzung-evaluation} berichteten
Kennzahlen; ihr jeweiliges Snapshot-Datum ist ausgewiesen. Der Snapshot der
produktiven Marktdatenbank sowie die redigierte Belegkette der
Echtgeld-Orders werden nach Abschluss der Echtgeld-Woche mit der nächsten
Archivversion ergänzt (kursive Einträge).

%% ---------------------------------------------------------------------------
%% GENERIERT durch build_quellcode_archiv.py am 27.07.2026 12:00 — NICHT von Hand
%% pflegen; Aenderungen in archiv_manifest.csv vornehmen und Skript ausfuehren.
%% ---------------------------------------------------------------------------
\begin{longtable}{@{}p{0.28\textwidth}p{0.17\textwidth}p{0.45\textwidth}@{}}
  \caption[Datenst\"ande des Archivs]{Ordner \texttt{Daten/} des Archivs: Daten- und Evaluationsst\"ande mit Snapshot-Datum; kursiv gesetzte Eintr\"age folgen mit der n\"achsten Archivversion nach Abschluss der Echtgeld-Woche.\quelleeigen}\label{tab:anhang-archiv-daten}\\
  \toprule
  Datei & Datenstand & Inhalt \\
  \midrule
  \endfirsthead
  \multicolumn{3}{@{}l}{\itshape Fortsetzung von vorheriger Seite}\\
  \toprule
  Datei & Datenstand & Inhalt \\
  \midrule
  \endhead
  \midrule
  \multicolumn{3}{r@{}}{\itshape Fortsetzung auf der n\"achsten Seite}\\
  \endfoot
  \bottomrule
  \endlastfoot
  \texttt{backtest\_\allowbreak protokoll\_\allowbreak 20260723.txt} & 23.07.2026 & Backtest-Protokoll B0/B1/B2 für 2.000 EUR und 1.000 EUR Kapitaleinsatz, Grundlage der Kennzahlen in Kapitel 3.4 und im Anhang Evaluationsbelege \\
  \texttt{kostentabelle.csv} & 19.07.2026 & Ex-Ante-Transaktionskosten der Bankmessung (Status 201, bankseitig validiert) \\
  \texttt{ml\_\allowbreak walkforward.csv} & 22.07.2026 & Walk-Forward-Kalibrierung der statistischen Prüfinstanz (23.803 bewertete Prognosetage) \\
  \texttt{gspc\_\allowbreak bars1d.csv} & 24.07.2026 & Tagesbars S\&P 500 (Export der Tabelle bars\_1d der Marktdatenbank, ab 30.12.1927) \\
  \texttt{marketdata\_\allowbreak snapshot.sqlite} & \itshape folgt (Live-Woche ab 27.07.2026) & \itshape Snapshot der produktiven Marktdatenbank (collector-vm) zum Abschluss der Echtgeld-Woche \\
  \texttt{journal\_\allowbreak echtgeld\_\allowbreak woche/} & \itshape folgt (Live-Woche ab 27.07.2026) & \itshape Redigierte Journal- und Berichtsauszüge der Echtgeld-Woche (Belegkette der Orders) \\
  \texttt{regime\_\allowbreak profit\_\allowbreak ergebnis.csv} & 24.07.2026 & Ergebnistabelle der Regime-Profitabilitätsanalyse (smatrend/sma über Kapitalstufen, Fenster und Kostenannahmen) — Grundlage der Aussagen zur Handelsgröße \\
  \texttt{veto\_\allowbreak overlay\_\allowbreak ergebnis.csv} & 24.07.2026 & Ergebnistabelle des ML-Veto-Overlays auf den Regime-Arm (Befund negativ) \\
  \texttt{backtest\_\allowbreak v6.csv} & 24.07.2026 & Walk-Forward-Backtest der ML-Prüfinstanz, Konventionsstand v6 \\
  \texttt{mlforecast\_\allowbreak lab\_\allowbreak ergebnis.csv} & 24.07.2026 & Laborergebnisse der Horizont- und Cross-Index-Untersuchung (Zusammenfassung) \\
  \texttt{mlforecast\_\allowbreak lab\_\allowbreak voll.csv} & 24.07.2026 & Laborergebnisse der Horizont- und Cross-Index-Untersuchung (Vollausgabe) \\
  \texttt{lab5\_\allowbreak metriken.csv} & 24.07.2026 & Kennzahlen des Laborlaufs 5 \\
  \texttt{lab5\_\allowbreak tail.csv} & 24.07.2026 & Randverteilung des Laborlaufs 5 \\
  \texttt{lab6\_\allowbreak metriken.csv} & 24.07.2026 & Kennzahlen des Laborlaufs 6 \\
  \texttt{lab6\_\allowbreak tail.csv} & 24.07.2026 & Randverteilung des Laborlaufs 6 \\
  \texttt{orderbuch\_\allowbreak smoketest\_\allowbreak 20260724.csv} & 24.07.2026 & Orderbuchauszug der Orderpfad-Erprobung (Anlage, Änderung, Storno) — Beleg für Statuswerte und Ordermerkmale der Bank \\
  \texttt{venue\_\allowbreak probe2\_\allowbreak protokoll\_\allowbreak 20260724.txt} & 24.07.2026 & Protokoll der Handelsplatz- und Ordertypprüfung inkl. abgelehnter Varianten (HTTP 422) — Beleg für die Grenzen der Bank-API \\
  \texttt{dimensions2\_\allowbreak DE000SB295Z1.json} & 24.07.2026 & Handelsplatz- und Ordertyp-Dimensionen der Bank für das Faktor-Zertifikat — maschineller Beleg zum Protokoll der Handelsplatzprüfung \\
  \texttt{trailing\_\allowbreak validation\_\allowbreak DE000SB295Z1.json} & 24.07.2026 & Validierungsantworten der Bank zu Trailing-/Stop-Varianten inkl. abgelehnter Kombinationen (HTTP 422) \\
\end{longtable}


%  =============================================================================
%  Ende
%  =============================================================================

%  =============================================================================
%  Datei-Hinweise & Konfiguration: siehe config.md
%  -> Abschnitt "anhang/evaluationsbelege.tex"
%
%  Anhang "Evaluationsbelege": reproduzierbare Belege der in Kapitel 3.4
%  genannten Kennzahlen. Druckaufbereitung der Protokolle analog zu
%  anhang/quelltexte.tex (ASCII-Ersetzungstabelle, footnotesize).
%  =============================================================================

\chapter{Evaluationsbelege}
\label{sec:anhang-evaluationsbelege}

Dieser Anhang dokumentiert die Datengrundlagen der Evaluation
(\cref{sec:umsetzung-evaluation}). Sämtliche Protokolle wurden unverändert von
den Produktionssystemen übernommen und lediglich typografisch für den Druck
aufbereitet; Umfangreiche Rohdaten (Tagesbars, Rohticks) sind durch
Herkunftsangaben referenziert, statt abgedruckt.

\section{Backtest-Protokoll (B0/B1/B2)}
\label{sec:anhang-backtest}

Vollständige Ausgabe von \texttt{strategy.py backtest} für die Kapitaleinsätze
2.000\,€ und 1.000\,€ (Datenstand 22.07.2026, 2.520 Handelstage; erzeugt am
23.07.2026 auf der agent-vm). Die in \cref{sec:umsetzung-evaluation-ergebnisse}
genannten Kennzahlen entstammen unverändert diesem Protokoll.

\lstinputlisting[basicstyle=\ttfamily\footnotesize,breaklines=true]{anhang/belege/backtest_protokoll_20260723.txt}

\section{Walk-Forward-Kalibrierung der statistischen Prüfinstanz}
\label{sec:anhang-walkforward}

Kalibrierungstabelle der Walk-Forward-Prüfung (\texttt{ml\_walkforward.csv};
23.803 bewertete Prognosetage): je Wahrscheinlichkeits-Bin die Anzahl der
Prognosen, die mittlere prognostizierte Wahrscheinlichkeit und die beobachtete
Trefferrate.

\begin{table}[htbp]
  \centering
  \begin{tabular}{r r r r}
    \hline
    Bin & $n$ & $\bar{p}$ prognostiziert & beobachtete Rate \\
    \hline
    0{,}1 & 2      & 0{,}184 & 0{,}000 \\
    0{,}2 & 7      & 0{,}262 & 0{,}429 \\
    0{,}3 & 127    & 0{,}364 & 0{,}504 \\
    0{,}4 & 4.142  & 0{,}483 & 0{,}501 \\
    0{,}5 & 19.434 & 0{,}524 & 0{,}525 \\
    0{,}6 & 91     & 0{,}624 & 0{,}549 \\
    \hline
  \end{tabular}
  \caption[Walk-Forward-Kalibrierung der statistischen Prüfinstanz]{%
    Walk-Forward-Kalibrierung der statistischen Prüfinstanz
    (Quelle: \texttt{ml\_walkforward.csv}, Stand 22.07.2026).\quelleeigen}
  \label{tab:anhang-walkforward}
\end{table}

\section{Ex-Ante-Transaktionskosten (Bankmessung)}
\label{sec:anhang-kosten}

Rohdaten der Kostenmessung gegen das Produktivsystem der Bank
(\texttt{cost\_probe.py}, Ex-Ante-Kostenausweis der Ordervalidierung, Status
201 = bankseitig validiert; Handelsplatz Live Trading/Direkthandel):

\begin{table}[htbp]
  \centering
  \begin{tabular}{r r r r r r r}
    \hline
    Volumen & Stück & Provision & Produktkosten & Kauf gesamt & Round-Trip & RT-Quote \\
    \hline
    150\,€ & 10 & 3{,}90\,€ & 0{,}14\,€ & 4{,}04\,€ & 7{,}80\,€ & 5{,}20\,\% \\
    495\,€ & 33 & 3{,}90\,€ & 0{,}47\,€ & 4{,}37\,€ & 7{,}80\,€ & 1{,}58\,\% \\
    990\,€ & 66 & 3{,}90\,€ & 0{,}95\,€ & 4{,}85\,€ & 7{,}80\,€ & 0{,}79\,\% \\
    \hline
  \end{tabular}
  \caption[Ex-Ante-Transaktionskosten der Bankmessung]{%
    Ex-Ante-Transaktionskosten der Bankmessung
    (Quelle: \texttt{kostentabelle.csv}, erhoben 19.07.2026).\quelleeigen}
  \label{tab:anhang-kosten}
\end{table}

\section{Herkunftsnachweis der Marktdatenbasis}
\label{sec:anhang-provenienz}

Herkunfts- und Integritätsnachweis der produktiven Marktdatenbank
(\texttt{marketdata.sqlite}, collector-vm; Abfragen vom 23.07.2026). Die
Rohtick-Tabelle enthält konstruktionsbedingt nur die jüngsten Handelstage, da
ältere Rohticks nach zehn Tagen zu Bars verdichtet und beschnitten werden
(vgl. \cref{sec:umsetzung-implementierung}).

\begin{table}[htbp]
  \centering
  \begin{tabular}{l l}
    \hline
    Merkmal & Wert \\
    \hline
    Tagesbars \texttt{bars\_1d} (S\&P~500) & 30.12.1927 -- 22.07.2026, 24.755 Zeilen \\
    Rohticks \texttt{ticks\_raw} (rollierend) & 79.945 Zeilen \\
    Integrität (\texttt{md5}) & \texttt{24771c148ec062b2c2dfa790b9d22975} \\
    Replikation (Litestream) & \texttt{gs://fom-ki-marketdata-backup/marketdata/} \\
    \hline
  \end{tabular}
  \caption[Herkunftsnachweis der Marktdatenbasis]{Herkunftsnachweis der
    Marktdatenbasis (collector-vm, Stand 23.07.2026).\quelleeigen}
  \label{tab:anhang-provenienz}
\end{table}

Die jüngsten Zeilen der Zweitquellenprüfung (\texttt{source\_crosscheck};
Spalte \texttt{stooq\_close} enthält aus historischen Gründen den
FRED-Schlusskurs) belegen die laufende Übereinstimmung der Datenquellen:

\begin{table}[htbp]
  \centering
  \begin{tabular}{l r r r r}
    \hline
    Datum & Primärquelle & FRED & Divergenz & Suspekt \\
    \hline
    21.07.2026 & 7.509{,}20 & 7.509{,}20 & 0{,}0\,\% & 0 \\
    20.07.2026 & 7.443{,}28 & 7.443{,}28 & 0{,}0\,\% & 0 \\
    17.07.2026 & 7.457{,}69 & 7.457{,}69 & 0{,}0\,\% & 0 \\
    16.07.2026 & 7.533{,}77 & 7.533{,}77 & 0{,}0\,\% & 0 \\
    15.07.2026 & 7.572{,}40 & 7.572{,}40 & 0{,}0\,\% & 0 \\
    14.07.2026 & 7.543{,}59 & 7.543{,}59 & 0{,}0\,\% & 0 \\
    13.07.2026 & 7.515{,}34 & 7.515{,}34 & 0{,}0\,\% & 0 \\
    10.07.2026 & 7.575{,}39 & 7.575{,}39 & 0{,}0\,\% & 0 \\
    09.07.2026 & 7.543{,}64 & 7.543{,}64 & 0{,}0\,\% & 0 \\
    \hline
  \end{tabular}
  \caption[Zweitquellenprüfung der Tagesschlusskurse]{Jüngste Zeilen der
    Zweitquellenprüfung gegen \glsxtrshort{fred}
    (Stand 23.07.2026).\quelleeigen}
  \label{tab:anhang-crosscheck}
\end{table}

% ---------------------------------------------------------------------------
% PLATZHALTER (wird nach Go-live 27.07. befüllt):
%
% \section{Belegkette der ersten Echtgeld-Order}
%   -> Tagesreport (redigiert), Journal-Auszug,
%      Ordervalidierung Status 200/201 (BELEGKONZEPT, Abschnitt B3).
% ---------------------------------------------------------------------------

% =============================================================================
%  Datei-Hinweise & Konfiguration: siehe config.md
%  -> Abschnitt "anhang/ki-prompts.tex"
% =============================================================================

\chapter{Anhang: Verwendete KI-Prompts}
\label{cha:anhang-prompts}

Im Verlauf der Erstellung dieser Arbeit wurden vier generative KI-Werkzeuge herangezogen. Anthropic Claude (Opus~5) diente der Erstellung des Glossars, des Formelzeichenverzeichnisses und des Anhangs, dem Entwurf der Marktdatenerfassung und ihrer Persistenz, dem Abgleich der Schnittstellendokumentation der Bank sowie der Anpassung des Extraktionsskripts für die Code- und Datenquellen der Cloud-Umgebung, der Erzeugung der Testskripte zur Ergebniskontrolle und der Testsuite für die Betriebsskripte, die Ausführungspfade und die Mutationsprobe. Anthropic Claude (Fable~5) unterstützte bei den Python- und Shell-Skripten des Agenten, bei der Cloud-Run-Anwendung des öffentlichen Dashboards, beim Totmann-Schalter, bei der statistischen Prüfinstanz samt Prognosemodul, bei der Gegenprüfung der Testskripte sowie bei der Erzeugung der beiden als KI-generiert gekennzeichneten Abbildungen. Google Gemini (3.5~Flash) wurde für die Rechtschreib- und Interpunktionsprüfung, die Quellenrecherche, den Abgleich der automatischen Tests und die Erstellung des Abkürzungsverzeichnisses eingesetzt. Moonshot AI Kimi (K3) diente einem allgemeinen Codereview der beigelegten Quellcodebeilage und dem Erstentwurf des Änderungsverzeichnisses der Beilage.

Gemäß FOM-Leitfaden Kap.~6.2.4 sind die genutzten Prompts in Kurzform — ohne die jeweils generierten Antworten — zu dokumentieren. Sämtliche Antworten wurden inhaltlich geprüft, eigenständig verifiziert und nicht wörtlich in die Arbeit übernommen.

% \kiwerkzeug{Tool}{Version/Stand}{Zugriffs-URL} -- einheitliche Darstellung
% eines KI-Werkzeugs; das zugehoerige \nocite{<bibkey>} steht in der
% jeweiligen Tool-Datei (Eintrag im KI-Verzeichnis, Kap. 6.2.4).
\newcommand{\kiwerkzeug}[3]{%
    \paragraph*{Verwendetes Tool}
    \begin{description}[leftmargin=2.5cm,style=nextline,labelwidth=2cm]
        \item[Tool] #1
        \item[Version] #2
        \item[Zugriff] \url{#3}
    \end{description}%
}

% --- Eingesetzte KI-Werkzeuge (je Werkzeug eine modulare Datei) -------------
% REIHENFOLGE NICHT AENDERN: Sie entspricht der Sortierung des KI-Verzeichnisses
% (biblatex sortiert nach Urheber: Anthropic, Anthropic, Google DeepMind,
% Moonshot AI). Die \note-Felder der Bib-Eintraege verweisen mit Anhang III.I
% bis III.IV auf genau diese Abschnittsnummern.
%  =============================================================================
%  Datei-Hinweise & Konfiguration: siehe config.md
%  -> Abschnitt "anhang/ki-prompts/tool-claude.tex"
%  =============================================================================

\section{Anthropic Claude (Opus~5)}
\label{sec:prompts-claude}

\kiwerkzeug{Anthropic Claude, Modell Opus 5}{Modellbezeichner
\texttt{claude-opus-5}, veröffentlicht am 24.07.2026}{https://claude.ai}
% Eintrag im KI-Verzeichnis (Leitfaden Kap. 6.2.4) erzwingen mittels \nocite.
\nocite{claudeOpus2026}

Eingesetzt für die Verzeichnisse des Vorspanns und den Aufbau des Anhangs
sowie für zwei Werkzeuge der Projektumgebung: die Anpassung des
Extraktionsskripts, das Quellcode- und Datenbestände der Google-Cloud-Umgebung
für das elektronische Archiv aufbereitet, und die Erzeugung der Testskripte
zur Ergebniskontrolle.

\subsection{Erstellung des Glossars}
\label{sec:prompts-claude-glossar}
\begin{enumerate}
    \item \enquote{Formuliere zu den folgenden Fachbegriffen jeweils eine
            prägnante, allgemeinverständliche Glossardefinition von höchstens
            zwei Sätzen im sachlichen wissenschaftlichen Stil: [\dots]}
            \par\hfill\textit{Kontext: Glossar (Vorspann)}

    \item \enquote{Prüfe die folgenden Glossareinträge auf fachliche
            Korrektheit, Widerspruchsfreiheit und einheitliche Formulierung:
            [\dots]}
            \par\hfill\textit{Kontext: Glossar (Vorspann)}
\end{enumerate}

\subsection{Erstellung des Formelzeichenverzeichnisses}
\label{sec:prompts-claude-formelzeichen}
\begin{enumerate}
    \item \enquote{Extrahiere aus den folgenden nummerierten Gleichungen
            sämtliche Formelzeichen und formuliere je Zeichen eine
            Kurzerklärung, die Bedeutung und Wertebereich beziehungsweise
            Einheit benennt: [\dots]}
            \par\hfill\textit{Kontext: Formelzeichenverzeichnis (Vorspann)}

    \item \enquote{Prüfe, ob jedes in den Gleichungen verwendete Formelzeichen
            im Verzeichnis geführt wird und ob umgekehrt jeder Eintrag des
            Verzeichnisses im Text tatsächlich vorkommt; nenne Abweichungen in
            beide Richtungen: [\dots]}
            \par\hfill\textit{Kontext: Formelzeichenverzeichnis (Vorspann)}
\end{enumerate}

\subsection{Aufbau des Anhangs}
\label{sec:prompts-claude-anhang}
\begin{enumerate}
    \item \enquote{Schlage eine Gliederung für einen dreiteiligen Anhang vor,
            der ein elektronisches Quellcodearchiv, die Belege einer
            Evaluation und die verwendeten KI-Prompts umfasst; benenne je
            Teil, welche Angaben erforderlich sind, damit ein Dritter die
            Ergebnisse nachvollziehen kann: [\dots]}
            \par\hfill\textit{Kontext: Anhang~\ref{cha:anhang-archiv} bis
            Anhang~\ref{cha:anhang-prompts}}

    \item \enquote{Formuliere zu den nachfolgenden Rohprotokollen jeweils
            einen einleitenden Absatz, der Herkunft, Erzeugungszeitpunkt und
            Aussagegehalt des Belegs benennt, ohne die Zahlen selbst zu
            wiederholen: [\dots]}
            \par\hfill\textit{Kontext:
            Anhang~\ref{sec:anhang-evaluationsbelege}}
\end{enumerate}

\subsection{Anpassung des Extraktionsskripts für Code- und Datenquellen}
\label{sec:prompts-claude-extraktion}

Das Skript zur Aufbereitung des Quellcodearchivs lag aus vorangegangenen
Projekten bereits vor; Gegenstand der Prompts war ausschließlich seine
Anpassung an die Zielarchitektur dieser Arbeit (\cref{fig:zielarchitektur})
mit ihren zwei virtuellen Maschinen, Journaldateien und zeitgesteuerten
Aufträgen.

\begin{enumerate}
    \item \enquote{Erweitere das vorhandene Extraktionsskript so, dass es
            neben Python- und Shell-Quelltexten auch die
            \texttt{systemd}-Einheiten und Cron-Definitionen beider
            virtueller Maschinen erfasst und je Datei Herkunftsmaschine sowie
            Ausführungsintervall in einem Manifest festhält: [\dots]}
            \par\hfill\textit{Kontext:
            Anhang~\ref{sec:anhang-archiv-quellcode}}

    \item \enquote{Ergänze das Skript um eine Redaktionsstufe, die
            Zugangsdaten sowie Konto- und Ordernummern aus den Journaldateien
            entfernt, bevor diese in das Archiv aufgenommen werden, und die
            jede Ersetzung protokolliert: [\dots]}
            \par\hfill\textit{Kontext: Anhang~\ref{sec:anhang-archiv-daten}}

    \item \enquote{Stelle sicher, dass die erzeugte Archivdatei bei
            unverändertem Inhalt reproduzierbar dieselbe Prüfsumme liefert,
            und erläutere, welche Zeitstempel dafür festzusetzen sind:
            [\dots]}
            \par\hfill\textit{Kontext:
            Anhang~\ref{sec:anhang-archiv-identitaet}}
\end{enumerate}

\subsection{Erzeugung der Testskripte zur Ergebniskontrolle}
\label{sec:prompts-claude-tests}
\begin{enumerate}
    \item \enquote{Entwirf zu der nachfolgenden Auswertungsfunktion eine
            Testsuite mit gemockten Schnittstellen, die jede
            sicherheitsrelevante Verzweigung abdeckt -- insbesondere die
            Pfade, die zu einer Ablehnung oder zu einem Handelsstopp führen:
            [\dots]}
            \par\hfill\textit{Kontext: \cref{sec:umsetzung-implementierung}}

    \item \enquote{Formuliere Testfälle, die prüfen, ob die im Protokoll
            ausgewiesenen Kennzahlen aus denselben Eingangsdaten
            reproduzierbar hervorgehen, und benenne die Toleranz, innerhalb
            derer Gleitkommaabweichungen zulässig sind: [\dots]}
            \par\hfill\textit{Kontext: \cref{sec:umsetzung-evaluation}}
\end{enumerate}

%  =============================================================================
%  Ende
%  =============================================================================

%  =============================================================================
%  Datei-Hinweise & Konfiguration: siehe config.md
%  -> Abschnitt "anhang/ki-prompts/tool-fable.tex"
%  =============================================================================

\section{Anthropic Claude (Fable~5)}
\label{sec:prompts-fable}

\kiwerkzeug{Anthropic Claude, Modell Fable 5}{Modellbezeichner
\texttt{claude-fable-5}, veröffentlicht am 09.06.2026}{https://claude.ai}
% Eintrag im KI-Verzeichnis (Leitfaden Kap. 6.2.4) erzwingen mittels \nocite.
\nocite{claudeFable2026}

Eingesetzt zur Unterstützung bei der Implementierung: Python- und
Shell-Skripte des Agenten, die Cloud-Run-Anwendung des öffentlichen
Dashboards, die Gegenprüfung der zuvor erzeugten Testskripte sowie die beiden
als KI-generiert gekennzeichneten Abbildungen dieser Arbeit
(\cref{fig:zielarchitektur} und \cref{fig:projektstruktur}).

\subsection{Unterstützung bei Python- und Shell-Skripten}
\label{sec:prompts-fable-skripte}
\begin{enumerate}
    \item \enquote{Überarbeite die folgende Funktion so, dass ein
            Verbindungsabbruch des Datenstroms nicht zum Prozessabbruch führt,
            sondern erkannt, protokolliert und über einen Wächterprozess
            neu verbunden wird; erhalte dabei die bestehende Signatur:
            [\dots]}
            \par\hfill\textit{Kontext: \cref{sec:umsetzung-architektur}}

    \item \enquote{Schreibe ein Shell-Skript, das den Entscheidungszyklus
            startet, eine Prozesssperre je Konto setzt, überlappende Läufe
            verhindert und im Fehlerfall mit eindeutigem Rückgabewert endet:
            [\dots]}
            \par\hfill\textit{Kontext: \cref{sec:umsetzung-implementierung}}

    \item \enquote{Erkläre, welche Fehlerbilder verteilter Systeme der
            nachfolgende Ausführungspfad noch nicht abfängt, und schlage je
            Fall die konservativste Behandlung vor: [\dots]}
            \par\hfill\textit{Kontext: \cref{sec:umsetzung-implementierung}}
\end{enumerate}

\subsection{Öffentliches Dashboard (Cloud-Run-Anwendung)}
\label{sec:prompts-fable-dashboard}
\begin{enumerate}
    \item \enquote{Entwirf einen zustandslosen Dienst, der redigierte
            JSON-Feeds aus einem Objektspeicher liest und daraus eine
            Weboberfläche rendert, ohne dass die handelnden Maschinen einen
            eingehenden Port öffnen müssen: [\dots]}
            \par\hfill\textit{Kontext: \cref{sec:umsetzung-implementierung}}

    \item \enquote{Erzeuge die Oberfläche des Dashboards so, dass der
            Indexverlauf mit den Durchschnittslinien und den Regimewechseln
            dargestellt wird und je Entscheidung das aggregierte Votum mit
            aufklappbaren Einzelbegründungen erscheint: [\dots]}
            \par\hfill\textit{Kontext: \cref{sec:umsetzung-implementierung}}

    \item \enquote{Prüfe den folgenden Rendering-Pfad darauf, ob Inhalte aus
            dem Journal ungeprüft in das Dokument geschrieben werden, und
            schlage eine Ausgabekodierung vor: [\dots]}
            \par\hfill\textit{Kontext: \cref{sec:umsetzung-implementierung}}
\end{enumerate}

\subsection{Verifikation der erzeugten Testskripte}
\label{sec:prompts-fable-testverifikation}

Die Testskripte wurden mit Claude Opus~5 entworfen
(\cref{sec:prompts-claude-tests}); die nachfolgenden Prompts dienten allein
ihrer unabhängigen Gegenprüfung durch ein zweites Modell.

\begin{enumerate}
    \item \enquote{Prüfe die folgende Testsuite kritisch: Welche Verzweigung
            des Produktivcodes wird von keinem Testfall erreicht, und welcher
            Test würde auch dann bestehen, wenn die geprüfte Zusicherung
            entfiele? [\dots]}
            \par\hfill\textit{Kontext: \cref{sec:umsetzung-implementierung}}

    \item \enquote{Benenne Testfälle, die nur scheinbar prüfen, was ihr Name
            behauptet, und schlage jeweils eine Formulierung vor, die den
            tatsächlich geprüften Sachverhalt trifft: [\dots]}
            \par\hfill\textit{Kontext: Anhang~\ref{sec:anhang-archiv-tests}}
\end{enumerate}

\subsection{Erzeugung der KI-generierten Abbildungen}
\label{sec:prompts-fable-abbildungen}

Beide Abbildungen sind in der jeweiligen Bildunterschrift als KI-generiert
ausgewiesen (Leitfaden Kap.~5.7 in Verbindung mit Kap.~6.2.4). Inhaltliche
Grundlage waren ausschließlich der eigene Quellcodebestand und die eigene
Architekturplanung; die Modellantwort lieferte die grafische Umsetzung.

\begin{enumerate}
    \item \enquote{Erzeuge eine SVG-Grafik der Zielarchitektur des
            Handelsagenten auf Basis der bereits erfolgten Implementierungen
            sowie der geplanten Bestandteile; umgesetzte Elemente
            durchgezogen, geplante gestrichelt, mit Legende. Style die Grafik
            attributbasiert ohne \texttt{<style>}-Block, damit die
            Inkscape-Konvertierung verlustfrei arbeitet: [\dots]}
            \par\hfill\textit{Kontext: \cref{fig:zielarchitektur}}

    \item \enquote{Erzeuge aus der nachfolgenden ASCII-Repräsentation der
            Verzeichnis- und Dateistruktur eine SVG-Grafik, die die Struktur
            in einer monospaced Schriftart schwarz auf weiß rendert, die
            Ebenen-Header durch Fettsatz hervorhebt und je Datei die
            Herkunftsmaschine kennzeichnet: [\dots]}
            \par\hfill\textit{Kontext: \cref{fig:projektstruktur}}
\end{enumerate}

%  =============================================================================
%  Ende
%  =============================================================================

%  =============================================================================
%  Datei-Hinweise & Konfiguration: siehe config.md
%  -> Abschnitt "anhang/ki-prompts/tool-gemini.tex"
%  =============================================================================

\section{Google Gemini (Flash~3.5)}
\label{sec:prompts-gemini}

\kiwerkzeug{Google Gemini, Modell Flash 3.5}{Stand: Mai 2026}{https://gemini.google.com}
% Eintrag im KI-Verzeichnis (Leitfaden Kap. 6.2.4) erzwingen mittels \nocite.
\nocite{geminiFlash2026}

\subsection{Literaturrecherche}
\label{sec:prompts-gemini-literatur}
\begin{enumerate}
    \item \enquote{Nenne einschlägige, zitierfähige Quellen (Monografien,
            Normen, begutachtete Fachartikel) zur BPMN-2.0-Prozessmodellierung
            und gib jeweils vollständige bibliografische Angaben an: [\dots]}
            \par\hfill\textit{Kontext: \Cref{cha:grundlagen}}

    \item \enquote{Fasse den aktuellen Forschungsstand zur
            Geschäftsprozessmodellierung im ERP-Umfeld zusammen und benenne
            die zentralen Referenzwerke: [\dots]}
            \par\hfill\textit{Kontext: \Cref{cha:grundlagen}}
\end{enumerate}

\subsection{Erstellung des Abkürzungsverzeichnisses}
\label{sec:prompts-gemini-abkuerzungen}
\begin{enumerate}
    \item \enquote{Extrahiere alle Abkürzungen aus dem folgenden Text und
            erstelle ein alphabetisch sortiertes Abkürzungsverzeichnis mit
            jeweiliger Langform: [\dots]}
            \par\hfill\textit{Kontext: gesamte Arbeit}
    \item \enquote{Überprüfe das vorhandene Abkürzungsverzeichnis auf
            Vollständigkeit und Konsistenz und ergänze fehlende Einträge: [\dots]}
            \par\hfill\textit{Kontext: gesamte Arbeit}
\end{enumerate}

%  =============================================================================
%  Ende
%  =============================================================================

%  =============================================================================
%  Datei-Hinweise & Konfiguration: siehe config.md
%  -> Abschnitt "anhang/ki-prompts/tool-kimi.tex"
%
%  BESONDERHEIT DIESES ABSCHNITTS: Die Prompts wurden in ENGLISCHER Sprache
%  gestellt und sind hier in deutscher Uebersetzung wiedergegeben. Der Hinweis
%  steht sowohl im Vorspann des Abschnitts als auch an jedem einzelnen Prompt
%  ("aus dem Englischen uebersetzt") -- beides beibehalten, damit kein Prompt
%  fuer sich betrachtet als deutschsprachige Eingabe missverstanden wird.
%
%  ABSCHNITTSNUMMER: Dieser Abschnitt ist Anhang III.IV. Das \note-Feld des
%  Bib-Eintrags moonshotKimi2026 verweist darauf. Wird die Reihenfolge der
%  \input-Zeilen in anhang/ki-prompts.tex geaendert, ist der Verweis dort
%  mitzufuehren.
%  =============================================================================

\section{Moonshot AI Kimi (K3)}
\label{sec:prompts-kimi}

\kiwerkzeug{Moonshot AI Kimi, Modell K3}{Modellbezeichner
\texttt{kimi-k3}, Modellvorstellung vom 17.07.2026}{https://www.kimi.com}
% Eintrag im KI-Verzeichnis (Leitfaden Kap. 6.2.4) erzwingen mittels \nocite.
\nocite{moonshotKimi2026}

Eingesetzt für ein allgemeines Codereview der beigelegten Quellcodebeilage
(\cref{cha:anhang-archiv}) sowie für den Erstentwurf des
Änderungsverzeichnisses der Beilage. Das Modell erhielt ausschließlich das
entpackte Archiv und keinen Zugriff auf den Projektordner; es war damit die
einzige Prüfinstanz, die den Code in genau der Form vorfand, in der er einem
Dritten vorliegt.

Die Prompts dieses Abschnitts wurden in englischer Sprache gestellt und sind
nachfolgend sinngleich ins Deutsche übersetzt wiedergegeben; der Zusatz
\textit{aus dem Englischen übersetzt} weist an jedem Prompt darauf hin. Die
Übersetzung war nicht Teil der Eingabe an das Modell.

Wie bei den übrigen Werkzeugen wurde jede Antwort eigenständig gegengeprüft.
Das Ergebnis fiel hier geteilt aus und ist deshalb ausdrücklich vermerkt: Von
zwei vorgeschlagenen Änderungen wurde die erste nach eigener Messung
übernommen -- sie betraf den Selbstschutz des Sitzungserhalts und beschrieb
einen zuvor unentdeckten Fehler --, die zweite als Rückschritt verworfen. Ihr
lag eine Verwechslung zweier Spalten des Archivmanifests zugrunde, die ohne
Kenntnis des Bauprozesses nicht erkennbar war; ihre Übernahme hätte den Bau
der Beilage abgebrochen. Auch der Erstentwurf des Änderungsverzeichnisses
wurde vor der Übernahme in drei Punkten korrigiert.

\subsection{Allgemeines Codereview der Quellcodebeilage}
\label{sec:prompts-kimi-codereview}
\begin{enumerate}
    \item \enquote{Anbei der vollständige Quellcodebestand eines
            Handelsagenten in der Fassung, die einer wissenschaftlichen
            Arbeit als Beilage beigefügt wird. Prüfe ihn ohne Vorgabe eines
            Schwerpunkts auf Fehler und benenne je Befund die betroffene
            Datei, die Zeile, die Wirkung im Betrieb und deine Sicherheit über
            den Befund: [\dots]}
            \par\hfill\textit{Kontext: \cref{cha:anhang-archiv}; aus dem
            Englischen übersetzt}

    \item \enquote{Beschränke die Prüfung nun auf die Shell-Skripte und dort
            auf die Behandlung von Rückgabewerten: Welche Abfrage deutet einen
            Fehlerausgang als regulären Zustand, und in welchem Fall bleibt
            eine notwendige Handlung dadurch unbemerkt aus? [\dots]}
            \par\hfill\textit{Kontext: \cref{cha:anhang-archiv},
            \cref{sec:umsetzung-implementierung}; aus dem Englischen
            übersetzt}

    \item \enquote{Lege zu jedem übernahmewürdigen Befund eine minimale
            Änderung als Textunterschied vor, die die bestehende Struktur
            wahrt, und benenne für jede Änderung, in welchen Umgebungen sie
            das Verhalten ändert und in welchen sie nachweislich wirkungslos
            bleibt: [\dots]}
            \par\hfill\textit{Kontext: \cref{cha:anhang-archiv}; aus dem
            Englischen übersetzt}
\end{enumerate}

\subsection{Entwurf des Änderungsverzeichnisses der Beilage}
\label{sec:prompts-kimi-changelog}
\begin{enumerate}
    \item \enquote{Formuliere zu den bestätigten Befunden einen Abschnitt für
            den Anhang, der die Unterschiede zwischen der vorliegenden und der
            nächsten Version der Beilage benennt: je Punkt das fehlerhafte
            Verhalten davor, die Änderung und den Beleg, dass sie wirkt --
            sachlich, ohne Wertung und ohne Wiedergabe von Quelltext:
            [\dots]}
            \par\hfill\textit{Kontext:
            Anhang~\ref{sec:anhang-archiv-identitaet}; aus dem Englischen
            übersetzt}

    \item \enquote{Prüfe den Abschnitt gegen die Angaben der Beilage: Welche
            Zahl, welcher Dateiname und welcher Verweis darin sind aus dem
            beigefügten Bestand nicht belegbar? [\dots]}
            \par\hfill\textit{Kontext:
            Anhang~\ref{sec:anhang-archiv-identitaet}; aus dem Englischen
            übersetzt}
\end{enumerate}

%  =============================================================================
%  Ende
%  =============================================================================


%  =============================================================================
%  Ende
%  =============================================================================


%  =============================================================================
%  Ende
%  =============================================================================


% Eigenstaendigkeitserklaerung (Leitfaden Kap. 7.2): "als letztes Blatt
% der Arbeit ... erhaelt jedoch keine Seitennummer und sie bleibt im
% Inhaltsverzeichnis unerwaehnt". Als formales Beidokument im
% anhang/-Ordner abgelegt.
% =============================================================================
%  Datei-Hinweise & Konfiguration: siehe config.md
%  -> Abschnitt "anhang/eigenstaendigkeitserklaerung.tex"
% =============================================================================

\cleardoublepage
\thispagestyle{empty}
\pagenumbering{gobble}

\section*{Eigenständigkeitserklärung}
\noindent
\textit{(schriftliche und mündliche Prüfungsleistungen; ausgenommen Klausuren)}

\bigskip

Hiermit versichere ich, dass ich die angemeldete Prüfungsleistung in allen
Teilen eigenständig ohne Hilfe von Dritten anfertigen und keine anderen als
die in der Prüfungsleistung angegebenen Quellen und zugelassenen Hilfsmittel
verwenden werde. Sämtliche wörtlichen und sinngemäßen Übernahmen inklusive
KI-generierter Inhalte werde ich kenntlich machen.

\medskip

Diese Prüfungsleistung hat zum Zeitpunkt der Abgabe weder in gleicher noch in
ähnlicher Form, auch nicht auszugsweise, bereits einer Prüfungsbehörde zur
Prüfung vorgelegen; hiervon ausgenommen sind Prüfungsleistungen, für die in
der Modulbeschreibung ausdrücklich andere Regelungen festgelegt sind.

\medskip

Mir ist bekannt, dass die Zuwiderhandlung gegen den Inhalt dieser Erklärung
einen Täuschungsversuch darstellt, der das Nichtbestehen der Prüfung zur Folge
hat und daneben strafrechtlich gem.~\S\,156 StGB verfolgt werden kann. Darüber
hinaus ist mir bekannt, dass ich bei schwerwiegender Täuschung exmatrikuliert
und mit einer Geldbuße bis zu 50.000~EUR nach der für mich gültigen
Rahmenprüfungsordnung belegt werden kann.

\medskip

Ich erkläre mich damit einverstanden, dass diese Prüfungsleistung zwecks
Plagiatsprüfung auf die Server externer Anbieter hochgeladen werden darf. Die
Plagiatsprüfung stellt keine Zurverfügungstellung für die Öffentlichkeit dar.

\vspace{2.5cm}

% Zwei Unterschriftsbloecke (Gruppenarbeit): je Autor Ort/Datum + Unterschrift.
\noindent
\begin{tabular}{@{}p{6cm}@{\hspace{1cm}}p{6cm}@{}}
\arbeitOrt, \abgabedatum & \\[0.9cm]
\hrulefill & \hrulefill \\
Ort, Datum & Unterschrift (\autorVorname{} \autorNachname) \\
\end{tabular}

\vspace{1.5cm}

\noindent
\begin{tabular}{@{}p{6cm}@{\hspace{1cm}}p{6cm}@{}}
\arbeitOrt, \abgabedatum &
    \raisebox{-0.1cm}{%
        \includesvg[height=1.1cm]{unterschrift_rolland}%
    } \\[-0.2cm]
\hrulefill & \hrulefill \\
Ort, Datum & Unterschrift (\autorZweiVorname{} \autorZweiNachname) \\
\end{tabular}

%  =============================================================================
%  Ende
%  =============================================================================


\end{document}

%  =============================================================================
%  Ende
%  =============================================================================
